\documentclass[12pt]{article}

\usepackage{amsmath, amssymb, amsthm}
\usepackage{scalerel,stackengine}
\usepackage{esint}
\usepackage{geometry}
\geometry{a4paper, margin=1in}
\usepackage{setspace}
\doublespacing
\usepackage{titlesec}
\usepackage[autostyle=true]{csquotes}
\usepackage{graphicx}
\usepackage{stmaryrd}
\usepackage{tikz}
\usetikzlibrary{arrows.meta, automata, positioning, trees}
\usepackage{fancyhdr}
\usepackage{dsfont}
\usepackage{imakeidx}
\usepackage{microtype}
\usepackage[svgnames]{xcolor}

\usepackage{algorithm}
\usepackage{algpseudocode}
\usepackage{listings}
\lstset{
    basicstyle=\ttfamily\small,
    keywordstyle=\color{NavyBlue}\bfseries,
    commentstyle=\color{Teal}\itshape,
    stringstyle=\color{DarkRed},
    frame=single,
    breaklines=true,
    numbers=left,
    numberstyle=\tiny\color{gray}
}

\usepackage{hyperref}
\hypersetup{
    colorlinks=true,
    linkcolor=NavyBlue,
    citecolor=DarkRed,
    urlcolor=Teal,
}

\usepackage[type={CC}, modifier={by-nc-sa}, version={4.0}]{doclicense}

\usepackage{enumitem}
\setlist{nosep}

\titleformat{\section}
{\normalfont\bfseries\LARGE}
{\thesection}{1em}{}
\titleformat{\subsection}
{\normalfont\bfseries\Large}
{\thesubsection}{1em}{}
\titleformat{\subsubsection}
{\normalfont\itshape\bfseries\large}
{\thesubsubsection}{1em}{}

\newtheoremstyle{blueplain}
  {3pt}{3pt}
  {\itshape}
  {}
  {\color{NavyBlue}\bfseries}
  {.}
  {0.5em}
  {}

\newtheoremstyle{bluedef}
  {3pt}{3pt}
  {\normalfont}
  {}
  {\color{Teal}\bfseries}
  {.}
  {0.5em}
  {}

\newtheoremstyle{blueremark}
  {3pt}{3pt}
  {\normalfont}
  {}
  {\color{DarkRed}\itshape}
  {.}
  {0.5em}
  {}

\theoremstyle{blueplain}
\newtheorem{theorem}{THEOREM}[section]
\newtheorem{proposition}[theorem]{Proposition}
\newtheorem{corollary}[theorem]{Corollary}
\newtheorem{lemma}[theorem]{Lemma}
\newtheorem{property}[theorem]{Property}

\theoremstyle{bluedef}
\newtheorem{definition}{DEFINITION}[section]
\newtheorem{problem}{Problem}[section]

\theoremstyle{blueremark}
\newtheorem{remark}{Remark}[section]
\newtheorem{example}{Example}[section]
\newtheorem{note}{Note}[section]

\newcommand{\R}{\mathbb{R}}
\newcommand{\N}{\mathbb{N}}
\newcommand{\C}{\mathbb{C}}
\newcommand{\Z}{\mathbb{Z}}
\newcommand{\Q}{\mathbb{Q}}
\newcommand{\eps}{\varepsilon}
\newcommand{\ups}{\upsilon}

\newcommand{\BigO}{\mathcal{O}}
\newcommand{\BigTheta}{\Theta}
\newcommand{\BigOmega}{\Omega}
\newcommand{\Error}{\textbf{Error}}

\makeindex[columns=2, title=Index, intoc]

\pagestyle{fancy}
\fancyhf{}
\fancyhead[L]{\textsc{\small Data Structures and Algorithm Analysis}} 
\fancyhead[R]{\small \nouppercase{\leftmark}}
\fancyfoot[R]{\hyperref[mytoc]{\Large $\hookleftarrow$}}
\fancyfoot[C]{\thepage}
\fancyfoot[L]{\href{https://creativecommons.org/licenses/by-nc-sa/4.0/deed.en}{\doclicenseIcon}}
\setlength{\headheight}{14pt}
\setlength{\headsep}{20pt}
\setlength{\footskip}{25pt}

\begin{document}

\pagenumbering{gobble}

\begin{titlepage}
    \centering
    \vspace*{3cm}

    {\Huge \bfseries Data Structures \\ and \\ Algorithm Analysis \par}

    \vspace{4cm}

    {\Large Kai Chen \par}
    \vspace{0.5cm}
    {\large \today \par}

    \vfill

    \rule{0.8\textwidth}{0.4pt} \\
    \vspace{0.5cm}
    {\small Part of the \href{https://github.com/kaiiiichen/SUSTech-Kai-Notes/wiki}{SUSTech-Kai-Notes} Initiative}
\end{titlepage}

\newpage
\thispagestyle{empty}

\vspace*{\fill}

\begin{center}
    \copyright\ \the\year\ Kai Chen.
    All rights reserved.

    \vspace{0.5cm}

    \doclicenseThis

    \vspace{1cm}

    This document was typeset using \href{https://www.latex-project.org}{\LaTeX}.
\end{center}
\clearpage

\pagenumbering{roman}

\section*{Preface}
These notes are compatible with the \href{https://github.com/kaiiiichen/SUSTech-Kai-Notes/tree/main/CS217%20Data%20Structures%20and%20Algorithm%20Analysis%20(H)}{CS217 Data Structures and Algorithm Analysis (H) (Fall 2025)} at \href{https://www.sustech.edu.cn/}{SUSTech}, and part of the course-notes-and-resources initiative: \href{https://github.com/kaiiiichen/SUSTech-Kai-Notes/wiki}{SUSTech-Kai-Notes}.

\medskip

\noindent Specific focus is placed on the mathematical analysis of algorithms, data structure implementation details, and complexity theory.

\clearpage

\phantomsection
\label{mytoc}

\onehalfspacing
\tableofcontents
\clearpage

\doublespacing
\pagenumbering{arabic}


\section{Getting Started}

\subsection{Foundations of Algorithms}

\subsubsection{Definitions}

\begin{definition}[Algorithm]\index{Algorithm}
    An \textbf{algorithm} is a well-defined computational procedure that takes some value, or set of values, as \textbf{input} and produces some value, or set of values, as \textbf{output}. It acts as a tool for solving a well-specified computational problem.
\end{definition}

\begin{problem}[The Sorting Problem]\index{Sorting Problem}
\textbf{Input}: A sequence of $n$ numbers $\langle a_{1},a_{2},...,a_{n}\rangle$.

\noindent \textbf{Output}: A permutation (reordering) $\langle{a_{1}}^{\prime},{a_{2}}^{\prime},...,{a_{n}}^{\prime}\rangle$ of the input sequence such that ${a_{1}}^{\prime}\le{a_{2}}^{\prime}\le\cdot\cdot\cdot\le{a_{n}}$.
\end{problem}

\begin{definition}[Instance]\index{Instance}
    An \textbf{instance} of a problem consists of the input (satisfying whatever constraints are imposed in the problem statement) needed to compute a solution.
    For example, the sequence $(31, 41, 59, 26, 41, 58)$ is an instance of the sorting problem.
\end{definition}

\begin{remark}
    We describe algorithms using \textbf{pseudocode}. This serves two purposes:
    \begin{enumerate}
        \item To show that algorithms exist independently of any particular programming language.
        \item To focus on ideas and logic rather than syntax issues or error-handling.
    \end{enumerate}
\end{remark}

\subsubsection{Correctness}

\begin{definition}[Correctness]\index{Correctness}
    An algorithm is \textbf{correct} if, for every input instance, it halts with the correct output. A correct algorithm solves the problem.
\end{definition}

\noindent Ideally, algorithms should be accompanied by a proof of correctness, rather than just testing on a few instances.

\subsection{The Computational Model}

To analyze the running time of algorithms effectively, we need a model that abstracts away specific hardware details (like clock rate, memory hierarchy, or multi-core architecture).

\begin{definition}[RAM Model]\index{RAM Model}
    The \textbf{Random-Access Machine (RAM)} model is a generic model of computation where instructions are executed one after another (no concurrent operations).
\end{definition}

\begin{property}[Elementary Operations]\index{Elementary Operations}
    The RAM model assumes that each elementary operation takes the same amount of time (a constant independent of the problem size). These operations include:
    \begin{itemize}
        \item \textbf{Arithmetic}: Add, subtract, multiply, divide, remainder, floor, ceiling.
        \item \textbf{Data Movement}: Variable assignments (storing, retrieving), copy.
        \item \textbf{Logical}: Comparisons, shifts, logical operations.
        \item \textbf{Control}: Conditional and unconditional branches (loops), subroutine calls and returns.
    \end{itemize}
\end{property}

\subsection{Insertion Sort}

\subsubsection{The Concept}

\textbf{Insertion Sort} works the way many people sort a hand of playing cards.
\begin{itemize}
    \item We start with an empty left hand and the cards face down on the table.
    \item We remove one card at a time from the table and insert it into the correct position in the left hand.
    \item To find the correct position for a card, we compare it with each of the cards already in the hand, from right to left.
    \item At all times, the cards held in the left hand are sorted.
\end{itemize}

\subsubsection{The Algorithm}

\begin{algorithm}
    \caption{Insertion Sort}\label{alg:insertionSort}
    \begin{algorithmic}[1]
        \Procedure{InsertionSort}{$A$}
        \For{$j = 2$ \textbf{to} $A.\text{length}$} \Comment{Iterate from the second element}
        \State $key = A[j]$
        \State \Comment{Insert $A[j]$ into the sorted sequence $A[1 \dots j-1]$}
        \State $i = j - 1$
        \While{$i > 0$ \textbf{and} $A[i] > key$}
        \State $A[i+1] = A[i]$ \Comment{Shift element to the right}
        \State $i = i - 1$
        \EndWhile
        \State $A[i+1] = key$ \Comment{Insert key into correct position}
        \EndFor
        \EndProcedure
    \end{algorithmic}
\end{algorithm}


\subsection{Correctness via Loop Invariants}

To prove that an algorithm (especially one with loops) is correct, we use the technique of \textbf{Loop Invariants}.

\begin{definition}[Loop Invariant]\index{Loop Invariant}
    A loop invariant is a statement that is always true and reflects the progress of the algorithm towards producing a correct output.
\end{definition}

To use a loop invariant to prove correctness, we must show three things:
\begin{enumerate}
    \item \textbf{Initialization}: The invariant is true before the first iteration of the loop.
    \item \textbf{Maintenance}: If the invariant is true before an iteration of the loop, it remains true before the next iteration.
    \item \textbf{Termination}: When the loop terminates, the invariant provides a useful property that helps show that the algorithm is correct.
\end{enumerate}

\subsubsection{Proof of Insertion Sort}

\begin{theorem}
    Algorithm \ref{alg:insertionSort} correctly sorts the array $A$.
\end{theorem}

\begin{proof}
    We use the following loop invariant:
    \begin{quote}
        \textit{"At the start of each iteration of the \textbf{for} loop (lines 1-8), the subarray $A[1 \dots j-1]$ consists of the elements originally in $A[1 \dots j-1]$, but in sorted order."}
    \end{quote}

    \noindent \textbf{Initialization}:
    We start with $j=2$. The subarray $A[1 \dots j-1]$ is $A[1 \dots 1]$, which consists of the single element $A[1]$. A single element is trivially sorted. Thus, the invariant holds.

    \noindent \textbf{Maintenance}:
    The body of the \textbf{for} loop works by moving $A[j-1], A[j-2], \dots$ one position to the right until the proper position for $A[j]$ (stored in $key$) is found. The $key$ is then placed into this position.
    At this point, the subarray $A[1 \dots j]$ consists of the elements originally in $A[1 \dots j]$, but in sorted order. Incrementing $j$ for the next iteration preserves the invariant.

    \noindent \textbf{Termination}:
    The loop terminates when $j = n + 1$. Substituting $n+1$ for $j$ in the invariant, we have that the subarray $A[1 \dots n]$ consists of the elements originally in $A[1 \dots n]$, but in sorted order. Since $A[1 \dots n]$ is the entire array, the algorithm is correct.
\end{proof}

\subsection{Runtime Analysis}

We analyze the running time $T(n)$ by summing the cost of each statement multiplied by the number of times it is executed.
\begin{itemize}
    \item Let $c_k$ be the cost of executing line $k$.
    \item Let $n$ be the size of the array $A$.
    \item Let $t_j$ be the number of times the \textbf{while} loop test is executed for a given value of $j$.
\end{itemize}

\subsubsection{Detailed Cost Breakdown}

\begin{itemize}
    \item Line 1 (for loop header): Executed $n$ times (checking $j$ from $2$ to $n+1$). Cost: $c_1 n$.
    \item Line 2 ($key = A[j]$): Executed $n-1$ times. Cost: $c_2 (n-1)$.
    \item Line 4 ($i = j-1$): Executed $n-1$ times. Cost: $c_4 (n-1)$.
    \item Line 5 (while test): Executed $\sum_{j=2}^n t_j$ times. Cost: $c_5 \sum_{j=2}^n t_j$.
    \item Line 6 (loop body, shift): Executed $\sum_{j=2}^n (t_j - 1)$ times. Cost: $c_6 \sum_{j=2}^n (t_j - 1)$.
    \item Line 7 (loop body, decrement): Executed $\sum_{j=2}^n (t_j - 1)$ times. Cost: $c_7 \sum_{j=2}^n (t_j - 1)$.
    \item Line 8 (assignment): Executed $n-1$ times. Cost: $c_8 (n-1)$.
\end{itemize}

The total running time $T(n)$ is:
\[
    T(n) = c_1 n + c_2 (n-1) + c_4 (n-1) + c_5 \sum_{j=2}^n t_j + c_6 \sum_{j=2}^n (t_j - 1) + c_7 \sum_{j=2}^n (t_j - 1) + c_8 (n-1)
\]

\subsubsection{Best Case Analysis}
The \textbf{best case} occurs when the array is already sorted.
In this scenario, $A[i] \le key$ immediately for every $j$. Thus, the \textbf{while} loop test is executed only once ($t_j = 1$), and the loop body is never executed.

Substituting $t_j = 1$:
\[ \sum_{j=2}^n t_j = \sum_{j=2}^n 1 = n-1 \]
\[ \sum_{j=2}^n (t_j - 1) = 0 \]

The runtime becomes:
\[
    T(n) = c_1 n + c_2 (n-1) + c_4 (n-1) + c_5 (n-1) + c_8 (n-1)
\]
\[
    T(n) = (c_1 + c_2 + c_4 + c_5 + c_8)n - (c_2 + c_4 + c_5 + c_8)
\]
This is a \textbf{linear function} of $n$: $T(n) = an + b = \Theta(n)$.

\subsubsection{Worst Case Analysis}
The \textbf{worst case} occurs when the array is reverse sorted.
In this scenario, we must compare the key with every element in the sorted subarray $A[1 \dots j-1]$. Thus, $t_j = j$.

We use the following summation formulas:
\[ \sum_{j=2}^n j = \frac{n(n+1)}{2} - 1 \]
\[ \sum_{j=2}^n (j-1) = \frac{n(n-1)}{2} \]

Substituting these into the total runtime equation:
\[
    T(n) = c_1 n + c_2(n-1) + c_4(n-1) + c_5\left(\frac{n(n+1)}{2} - 1\right) + (c_6 + c_7)\left(\frac{n(n-1)}{2}\right) + c_8(n-1)
\]
Grouping terms by powers of $n$:
\[
    T(n) = \left(\frac{c_5}{2} + \frac{c_6}{2} + \frac{c_7}{2}\right)n^2 + \left(c_1 + c_2 + c_4 + \frac{c_5}{2} - \frac{c_6}{2} - \frac{c_7}{2} + c_8\right)n - (c_2 + c_4 + c_5 + c_8)
\]
This is a \textbf{quadratic function} of $n$: $T(n) = an^2 + bn + c = \Theta(n^2)$.

\clearpage


\section{Runtime and Asymptotic Notation}

\subsection{Recap: Runtime of Insertion Sort}

In the previous lecture, we derived the running time function $T(n)$ for Insertion Sort on an input of size $n$. The exact formula depended on specific costs $c_i$ for each machine instruction:
\[
    T(n) = c_1 n + c_2(n-1) + c_4(n-1) + c_5 \sum_{j=2}^n t_j + \dots
\]
This level of detail is often too messy and hardware-dependent. We observed two extreme cases:
\begin{itemize}
    \item \textbf{Best Case} (Sorted Array): $T(n) = an + b$ (Linear).
    \item \textbf{Worst Case} (Reverse Sorted Array): $T(n) = an^2 + bn + c$ (Quadratic).
\end{itemize}

\subsubsection{Why Focus on the Worst Case?}
\begin{enumerate}
    \item \textbf{Guarantee}: The worst-case runtime gives us an absolute upper bound. The algorithm will never take longer than this.
    \item \textbf{Frequency}: For some algorithms, the worst case occurs fairly often (e.g., searching for a non-existent item in a database).
    \item \textbf{Average Case}: Often, the "average" case is roughly as bad as the worst case. For Insertion Sort, the average case is also quadratic.
\end{enumerate}

\subsection{Asymptotic Analysis}

\subsubsection{The Philosophy of Asymptotics}
To compare algorithms effectively, we need a metric that is independent of:
\begin{itemize}
    \item The specific machine (CPU speed, memory architecture).
    \item The programming language or compiler efficiency.
    \item Small input sizes (where most algorithms are fast enough).
\end{itemize}

\textbf{Key Idea}: We focus on the \textbf{rate of growth} of the running time as the input size $n$ approaches infinity ($n \to \infty$).
\begin{itemize}
    \item We ignore \textbf{lower-order terms}: For large $n$, the highest power dominates (e.g., in $n^2 + 100n$, the $n^2$ term is overwhelmingly larger).
    \item We ignore \textbf{leading constants}: Constants depend on hardware. We want to know if the algorithm scales linearly, quadratically, etc.
\end{itemize}

\subsection{Asymptotic Notations}

We use five standard notations to describe the asymptotic behavior of functions. Let $f(n)$ be the algorithm's runtime and $g(n)$ be a reference function.

\subsubsection{Big-Theta Notation ($\Theta$): Tight Bound}

\begin{definition}[Big-Theta $\Theta$]\index{Big-Theta}
    For a given function $g(n)$, $\Theta(g(n))$ is the set of functions:
    \[
        \Theta(g(n)) = \{ f(n) : \exists c_1, c_2 > 0 \text{ and } n_0 \ge 0 \text{ such that } 0 \le c_1 g(n) \le f(n) \le c_2 g(n) \text{ for all } n \ge n_0 \}
    \]
\end{definition}
\noindent \textbf{Intuition}: For sufficiently large $n$, $f(n)$ is "sandwiched" between two constant multiples of $g(n)$.
\begin{itemize}
    \item We say $f(n)$ grows \textbf{at the same rate} as $g(n)$.
    \item Example: $2n^2 + 3n + 1 = \Theta(n^2)$.
    \item Note: We often write $f(n) = \Theta(g(n))$ (equality) to mean $f(n) \in \Theta(g(n))$ (set membership).
\end{itemize}



\subsubsection{Big-O Notation ($O$): Upper Bound}

\begin{definition}[Big-O $O$]\index{Big-O}
    For a given function $g(n)$, $O(g(n))$ is the set of functions:
    \[
        O(g(n)) = \{ f(n) : \exists c > 0 \text{ and } n_0 \ge 0 \text{ such that } 0 \le f(n) \le c g(n) \text{ for all } n \ge n_0 \}
    \]
\end{definition}
\noindent \textbf{Intuition}: $f(n)$ grows \textbf{no faster than} $g(n)$. It gives an asymptotic \textbf{ceiling}.
\begin{itemize}
    \item Note: $f(n) = \Theta(g(n)) \implies f(n) = O(g(n))$.
    \item Ideally, we want the tightest upper bound, but loose bounds are technically correct (e.g., $2n = O(n^2)$ is true, but weak).
\end{itemize}



\subsubsection{Big-Omega Notation ($\Omega$): Lower Bound}

\begin{definition}[Big-Omega $\Omega$]\index{Big-Omega}
    For a given function $g(n)$, $\Omega(g(n))$ is the set of functions:
    \[
        \Omega(g(n)) = \{ f(n) : \exists c > 0 \text{ and } n_0 \ge 0 \text{ such that } 0 \le c g(n) \le f(n) \text{ for all } n \ge n_0 \}
    \]
\end{definition}
\noindent \textbf{Intuition}: $f(n)$ grows \textbf{at least as fast as} $g(n)$. It gives an asymptotic \textbf{floor}.



\begin{theorem}[Relationship Theorem]
    For any two functions $f(n)$ and $g(n)$:
    \[ f(n) = \Theta(g(n)) \iff f(n) = O(g(n)) \text{ AND } f(n) = \Omega(g(n)) \]
\end{theorem}

\subsubsection{Strict Bounds ($o$ and $\omega$)}

While $O$ and $\Omega$ correspond to $\le$ and $\ge$, the notations $o$ and $\omega$ correspond to $<$ and $>$.

\begin{definition}[Little-o $o$]\index{Little-o}
    $f(n) = o(g(n))$ if for \textbf{any} constant $c > 0$, there exists $n_0 > 0$ such that $0 \le f(n) < c g(n)$ for all $n \ge n_0$.
    Alternatively defined by limits:
    \[ \lim_{n \to \infty} \frac{f(n)}{g(n)} = 0 \]
    \textbf{Meaning}: $f(n)$ becomes insignificant relative to $g(n)$ as $n$ grows.
\end{definition}

\begin{definition}[Little-omega $\omega$]\index{Little-omega}
    $f(n) = \omega(g(n))$ if for \textbf{any} constant $c > 0$, there exists $n_0 > 0$ such that $0 \le c g(n) < f(n)$ for all $n \ge n_0$.
    Alternatively:
    \[ \lim_{n \to \infty} \frac{f(n)}{g(n)} = \infty \]
    \textbf{Meaning}: $f(n)$ grows strictly faster than $g(n)$.
\end{definition}

\subsection{Analogy with Real Numbers}

We can map the asymptotic relationships between functions to comparison operators for real numbers:

\begin{table}[h]
    \centering
    \begin{tabular}{|c|c|l|}
        \hline
        \textbf{Notation}     & \textbf{Analogy} & \textbf{Intuition}              \\ \hline
        $f(n) = O(g(n))$      & $f \le g$        & Upper bound (tight or loose)    \\ \hline
        $f(n) = \Omega(g(n))$ & $f \ge g$        & Lower bound (tight or loose)    \\ \hline
        $f(n) = \Theta(g(n))$ & $f = g$          & Tight bound (equal growth rate) \\ \hline
        $f(n) = o(g(n))$      & $f < g$          & Strictly smaller growth         \\ \hline
        $f(n) = \omega(g(n))$ & $f > g$          & Strictly larger growth          \\ \hline
    \end{tabular}
    \caption{Comparison of Asymptotic Notations}
\end{table}

\subsection{Properties and Growth Hierarchy}

\subsubsection{Properties}
\begin{itemize}
    \item \textbf{Transitivity}: If $f(n) = \Theta(g(n))$ and $g(n) = \Theta(h(n))$, then $f(n) = \Theta(h(n))$. (Holds for all 5 notations).
    \item \textbf{Reflexivity}: $f(n) = \Theta(f(n))$. (Holds for $O, \Omega$).
    \item \textbf{Symmetry}: $f(n) = \Theta(g(n)) \iff g(n) = \Theta(f(n))$.
    \item \textbf{Transpose Symmetry}: $f(n) = O(g(n)) \iff g(n) = \Omega(f(n))$.
\end{itemize}

\subsubsection{Common Growth Functions}
Ordered from slowest to fastest growth:
\begin{enumerate}
    \item $O(1)$: Constant time.
    \item $O(\lg n)$: Logarithmic time.
    \item $O(\sqrt{n})$: Square root.
    \item $O(n)$: Linear time.
    \item $O(n \lg n)$: Linearithmic (often seen in efficient sorting).
    \item $O(n^2)$: Quadratic.
    \item $O(n^3)$: Cubic.
    \item $O(2^n)$: Exponential.
    \item $O(n!)$: Factorial.
\end{enumerate}

\begin{remark}[Dominance Rules]
    \begin{itemize}
        \item \textbf{Polylogarithms vs Polynomials}: For any constants $a, b > 0$, $(\lg n)^a = o(n^b)$. (Any positive polynomial power beats any polylog).
        \item \textbf{Polynomials vs Exponentials}: For any constants $a > 0, b > 1$, $n^a = o(b^n)$. (Any exponential with base $>1$ beats any polynomial).
    \end{itemize}
\end{remark}

\subsection{Application to Insertion Sort}

Let $T(n)$ be the running time of Insertion Sort on an array of size $n$.

\begin{itemize}
    \item \textbf{Worst Case}: Occurs when the array is reverse sorted.
          \[ T(n) = an^2 + bn + c \implies T(n) = \Theta(n^2) \]
    \item \textbf{Best Case}: Occurs when the array is already sorted.
          \[ T(n) = an + b \implies T(n) = \Theta(n) \]
\end{itemize}

\begin{note}[Common Misconception]
    It is \textbf{incorrect} to say "The running time of Insertion Sort is $\Theta(n^2)$".
    Why? Because for the best-case input, the runtime is linear, not quadratic.
    \begin{itemize}
        \item \textbf{Correct Statement 1}: The \textit{worst-case} running time is $\Theta(n^2)$.
        \item \textbf{Correct Statement 2}: For \textit{any} input, the running time is $O(n^2)$ (upper bound) and $\Omega(n)$ (lower bound).
    \end{itemize}
\end{note}

\clearpage


\section{Divide-and-Conquer}

\subsection{The Paradigm Shift}

In previous lectures, we explored **Insertion Sort**, which follows an **Incremental Approach**: we process elements one by one and insert them into a growing sorted subarray.
To break the $\Theta(n^2)$ barrier, we need a fundamentally different strategy: **Divide-and-Conquer**.

\subsubsection{Divide-and-Conquer Strategy}
This paradigm involves three recursive steps:
\begin{enumerate}
    \item \textbf{Divide}: Break the problem into several smaller subproblems that are similar to the original problem but smaller in size.
    \item \textbf{Conquer}: Solve the subproblems recursively. If the subproblem sizes are small enough (the base case), solve them directly.
    \item \textbf{Combine}: Merge the solutions to the subproblems to create a solution to the original problem.
\end{enumerate}

\subsubsection{Motivating Example: Binary Search}
Consider finding a number $x$ in a \textbf{sorted} array of size $n$.
\begin{itemize}
    \item \textbf{Linear Search}: Scanning from start to end takes $\Theta(n)$ in the worst case.
    \item \textbf{Binary Search}: Check the middle element. If $x$ is smaller, discard the right half; otherwise, discard the left half.
          \[ T(n) = T(n/2) + \Theta(1) \implies T(n) = \Theta(\lg n) \]
          \textbf{Insight}: By dividing the problem size by a constant factor (2) at each step, we exponentially reduce the work, achieving logarithmic complexity.
\end{itemize}

\subsection{Merge Sort}

Merge Sort is the archetypal divide-and-conquer sorting algorithm.

\subsubsection{The Algorithm}
\begin{enumerate}
    \item \textbf{Divide}: Split the $n$-element sequence into two subsequences of size $n/2$.
    \item \textbf{Conquer}: Recursively sort the two subsequences using Merge Sort.
    \item \textbf{Combine}: Merge the two sorted subsequences to produce the final sorted answer.
\end{enumerate}
\textbf{Base Case}: A sequence of length 1 is trivially sorted.

\subsubsection{The Merge Procedure}
The core logic lies in the \textbf{Combine} step. The procedure \texttt{Merge(A, p, q, r)} assumes that the subarrays $A[p \dots q]$ and $A[q+1 \dots r]$ are already sorted. It merges them into a single sorted subarray $A[p \dots r]$.

\textbf{Mechanism (Two Pointers)}:
We view the two subarrays as piles of cards face up. We compare the top cards of each pile, pick the smaller one, place it in the output pile, and repeat.
To avoid checking for empty piles constantly, we can use a **sentinel value** ($\infty$) at the end of each subarray.

\begin{algorithm}
    \caption{Merge Procedure}\label{alg:merge}
    \begin{algorithmic}[1]
        \Procedure{Merge}{$A, p, q, r$}
        \State $n_1 = q - p + 1$
        \State $n_2 = r - q$
        \State Let $L[1 \dots n_1+1]$ and $R[1 \dots n_2+1]$ be new arrays
        \For{$i = 1$ \textbf{to} $n_1$}
        \State $L[i] = A[p + i - 1]$
        \EndFor
        \For{$j = 1$ \textbf{to} $n_2$}
        \State $R[j] = A[q + j]$
        \EndFor
        \State $L[n_1 + 1] = \infty$; $R[n_2 + 1] = \infty$ \Comment{Sentinels}
        \State $i = 1$; $j = 1$
        \For{$k = p$ \textbf{to} $r$}
        \If{$L[i] \le R[j]$}
        \State $A[k] = L[i]$
        \State $i = i + 1$
        \Else
        \State $A[k] = R[j]$
        \State $j = j + 1$
        \EndIf
        \EndFor
        \EndProcedure
    \end{algorithmic}
\end{algorithm}

\noindent \textbf{Runtime of Merge}:
The procedure performs a constant amount of work for each element in the subarrays. Specifically, the \textbf{for} loop runs $n = r - p + 1$ times.
\[ T_{\text{merge}}(n) = \Theta(n) \]


\subsubsection{Correctness of Merge}
We prove correctness using a **Loop Invariant**:
\begin{quote}
    \textit{At the start of each iteration of the \textbf{for} loop, the subarray $A[p \dots k-1]$ contains the $k-p$ smallest elements of $L$ and $R$, in sorted order. Furthermore, $L[i]$ and $R[j]$ are the smallest elements of their arrays that have not been copied back to $A$.}
\end{quote}
\begin{itemize}
    \item \textbf{Initialization}: For $k=p$, the subarray $A[p \dots p-1]$ is empty, which trivially satisfies the invariant.
    \item \textbf{Maintenance}: Suppose $L[i] \le R[j]$. Then $L[i]$ is the smallest uncopied element. We copy it to $A[k]$. Now $A[p \dots k]$ contains the $k-p+1$ smallest elements. Incrementing $k$ and $i$ maintains the invariant.
    \item \textbf{Termination}: When the loop ends, $k=r+1$. The subarray $A[p \dots r]$ contains all elements (except sentinels) in sorted order.
\end{itemize}

\subsection{Analysis of Merge Sort}

\subsubsection{Recurrence Relation}
Let $T(n)$ be the time to sort $n$ numbers.
\begin{itemize}
    \item \textbf{Divide}: Computing the middle index takes constant time $\Theta(1)$.
    \item \textbf{Conquer}: We recursively solve two subproblems, each of size $n/2$, contributing $2T(n/2)$.
    \item \textbf{Combine}: The \texttt{Merge} procedure takes time linear in the number of elements, $\Theta(n)$.
\end{itemize}
Thus, the recurrence is:
\[
    T(n) = \begin{cases}
        \Theta(1)           & \text{if } n=1   \\
        2T(n/2) + \Theta(n) & \text{if } n > 1
    \end{cases}
\]


\subsubsection{Recursion Tree Visualization}
To solve $T(n) = 2T(n/2) + cn$, visualize a tree:
\begin{itemize}
    \item \textbf{Top Level}: Cost $cn$.
    \item \textbf{Level 1}: Two subproblems of size $n/2$. Cost $2 \times c(n/2) = cn$.
    \item \textbf{Level 2}: Four subproblems of size $n/4$. Cost $4 \times c(n/4) = cn$.
    \item \dots
    \item \textbf{Depth}: The tree splits until $n=1$. The height is $\lg n$.
    \item \textbf{Total Cost}: There are $\lg n + 1$ levels, each costing $cn$.
          \[ \text{Total} \approx cn \times \lg n = \Theta(n \lg n) \]
\end{itemize}

\subsubsection{Comparison with Insertion Sort}
\begin{itemize}
    \item \textbf{Time}: Merge Sort is $\Theta(n \lg n)$ in the worst, best, and average cases. This is asymptotically faster than Insertion Sort's worst case $\Theta(n^2)$.
    \item \textbf{Space}: Merge Sort is \textbf{not} in-place. It requires $\Theta(n)$ auxiliary space for the `L` and `R` arrays in the `Merge` step. Insertion Sort is in-place ($O(1)$ space). This is the main drawback of Merge Sort.
\end{itemize}

\subsection{The Master Theorem}

The Master Theorem provides a "cookbook" method to solve recurrences of the form:
\[ T(n) = aT(n/b) + f(n) \]
where $a \ge 1$ (number of subproblems), $b > 1$ (shrinkage factor), and $f(n)$ (cost of divide/combine).

\subsubsection{The Watershed Function}
We compare the driving function $f(n)$ with the \textbf{watershed function} $n^{\log_b a}$.
\begin{itemize}
    \item $n^{\log_b a}$ represents the rate of growth of the leaves (the cost of the base cases).
    \item $f(n)$ represents the cost of the work done at the root (divide and combine).
\end{itemize}
It's a "tug-of-war" between the root work and the leaf work.

\subsubsection{The Three Cases}

\begin{theorem}[Master Theorem]\index{Master Theorem}
    \begin{enumerate}
        \item \textbf{Case 1 (Leaves Dominate)}: If $f(n) = O(n^{\log_b a - \epsilon})$ for some $\epsilon > 0$, then the recursion cost dominates:
              \[ T(n) = \Theta(n^{\log_b a}) \]
        \item \textbf{Case 2 (Balanced)}: If $f(n) = \Theta(n^{\log_b a} \lg^k n)$ (typically $k=0$), then the cost is distributed evenly across levels:
              \[ T(n) = \Theta(n^{\log_b a} \lg^{k+1} n) \]
        \item \textbf{Case 3 (Root Dominates)}: If $f(n) = \Omega(n^{\log_b a + \epsilon})$ for some $\epsilon > 0$, and if the regularity condition $af(n/b) \le cf(n)$ holds, then the root work dominates:
              \[ T(n) = \Theta(f(n)) \]
    \end{enumerate}
\end{theorem}

\subsubsection{Examples}

\textbf{Example 1: Merge Sort}
\[ T(n) = 2T(n/2) + n \]
\begin{itemize}
    \item $a=2, b=2, f(n)=n$.
    \item Watershed: $n^{\log_2 2} = n^1 = n$.
    \item Compare: $f(n) = \Theta(\text{Watershed})$. This is \textbf{Case 2} (with $k=0$).
    \item Result: $T(n) = \Theta(n \lg n)$.
\end{itemize}

\textbf{Example 2}
\[ T(n) = 9T(n/3) + n \]
\begin{itemize}
    \item $a=9, b=3, f(n)=n$.
    \item Watershed: $n^{\log_3 9} = n^2$.
    \item Compare: $f(n) = n = O(n^{2-\epsilon})$. Watershed dominates. This is \textbf{Case 1}.
    \item Result: $T(n) = \Theta(n^2)$.
\end{itemize}

\textbf{Example 3}
\[ T(n) = 3T(n/4) + n \lg n \]
\begin{itemize}
    \item $a=3, b=4, f(n)=n \lg n$.
    \item Watershed: $n^{\log_4 3} \approx n^{0.793}$.
    \item Compare: $f(n) = \Omega(n^{0.793+\epsilon})$. Root dominates. This is \textbf{Case 3}.
    \item Regularity: $3(n/4) \lg(n/4) \le c n \lg n$ holds for $c \approx 3/4$.
    \item Result: $T(n) = \Theta(n \lg n)$.
\end{itemize}

\clearpage


\section{Heapsort}

\subsection{Motivation: Smarter Selection}

To understand Heapsort, let's revisit \textbf{Selection Sort}.
\begin{itemize}
    \item \textbf{Selection Sort Strategy}: Repeatedly find the maximum element in the remaining unsorted array and move it to the correct position.
    \item \textbf{The Inefficiency}: Finding the maximum takes $\Theta(n)$ time because we scan the array linearly. We do this $n$ times, leading to $\Theta(n^2)$.
    \item \textbf{The Waste}: In every iteration, Selection Sort "forgets" all the comparisons it made. It starts from scratch.
\end{itemize}

\textbf{Heapsort's Insight}: Can we use a data structure to "memorize" the comparison results?
\begin{itemize}
    \item Instead of scanning linearly, we organize elements so that finding the maximum is trivial ($O(1)$).
    \item Removing the maximum and "repairing" the structure should be fast ($O(\lg n)$).
    \item This turns the "Rebuild" process (Selection Sort) into a "Repair" process (Heapsort).
\end{itemize}

\subsection{The Heap Data Structure}

\subsubsection{Definition}

A \textbf{Heap} (specifically, a binary heap) is an array object that we view as a \textbf{nearly complete binary tree}.
\begin{itemize}
    \item \textbf{Structure}: The tree is completely filled on all levels except possibly the lowest, which is filled from the left up to a point.
    \item \textbf{Implicit Representation}: We don't use pointers. The tree structure is implicit in the array indices.
\end{itemize}

For a node at index $i$:
\begin{itemize}
    \item $\text{Parent}(i) = \lfloor i/2 \rfloor$
    \item $\text{Left}(i) = 2i$
    \item $\text{Right}(i) = 2i + 1$
\end{itemize}

\subsubsection{The Heap Property}

There are two kinds of heaps:
\begin{enumerate}
    \item \textbf{Max-Heap}: For every node $i$ other than the root:
          \[ A[\text{Parent}(i)] \ge A[i] \]
          The largest element is stored at the root. (Used for Heapsort).
    \item \textbf{Min-Heap}: For every node $i$ other than the root:
          \[ A[\text{Parent}(i)] \le A[i] \]
          The smallest element is stored at the root. (Used for Priority Queues).
\end{enumerate}

\subsection{Core Operations}

\subsubsection{Maintaining the Property: Max-Heapify}

\textbf{Problem}: Suppose the binary trees rooted at $\text{Left}(i)$ and $\text{Right}(i)$ are max-heaps, but $A[i]$ might be smaller than its children, violating the max-heap property.
\textbf{Solution}: Let the value at $A[i]$ "float down" until the property is restored.

\begin{algorithm}
    \caption{Max-Heapify}\label{alg:maxHeapify}
    \begin{algorithmic}[1]
        \Procedure{Max-Heapify}{$A, i$}
        \State $l = \text{Left}(i)$
        \State $r = \text{Right}(i)$
        \State $largest = i$
        \If{$l \le A.\text{heap-size}$ \textbf{and} $A[l] > A[largest]$}
        \State $largest = l$
        \EndIf
        \If{$r \le A.\text{heap-size}$ \textbf{and} $A[r] > A[largest]$}
        \State $largest = r$
        \EndIf
        \If{$largest \ne i$}
        \State Exchange $A[i]$ with $A[largest]$
        \State \Call{Max-Heapify}{$A, largest$} \Comment{Recursive call on the affected subtree}
        \EndIf
        \EndProcedure
    \end{algorithmic}
\end{algorithm}

\noindent \textbf{Runtime}: The running time $T(n)$ depends on the height of the node $h$. In the worst case, we traverse from the root to a leaf.
\[ T(n) = O(h) = O(\lg n) \]

\subsubsection{Building the Heap: Build-Max-Heap}

We can convert an arbitrary array $A[1 \dots n]$ into a max-heap by calling \texttt{Max-Heapify} in a bottom-up manner.
\textbf{Observation}: The elements in the subarray $A[\lfloor n/2 \rfloor + 1 \dots n]$ are all leaves. Leaves are trivially max-heaps of size 1. We start from their parents and work up to the root.

\begin{algorithm}
    \caption{Build-Max-Heap}\label{alg:buildMaxHeap}
    \begin{algorithmic}[1]
        \Procedure{Build-Max-Heap}{$A$}
        \State $A.\text{heap-size} = A.\text{length}$
        \For{$i = \lfloor A.\text{length}/2 \rfloor$ \textbf{downto} $1$}
        \State \Call{Max-Heapify}{$A, i$}
        \EndFor
        \EndProcedure
    \end{algorithmic}
\end{algorithm}

\noindent \textbf{Runtime Analysis (Crucial)}:
A naive bound is $O(n \lg n)$ because we call \texttt{Max-Heapify} $O(n)$ times. However, this is not tight.
\begin{itemize}
    \item Most nodes have small heights.
    \item \texttt{Max-Heapify} takes time proportional to the height $h$ of the node.
    \item A heap of size $n$ has at most $\lceil n/2^{h+1} \rceil$ nodes of height $h$.
\end{itemize}
Total work:
\[
    \sum_{h=0}^{\lfloor \lg n \rfloor} \left\lceil \frac{n}{2^{h+1}} \right\rceil O(h) = O\left( n \sum_{h=0}^{\lfloor \lg n \rfloor} \frac{h}{2^h} \right)
\]
Using the summation formula $\sum_{h=0}^{\infty} \frac{h}{2^h} = 2$, we get:
\[ T(n) = O(n \cdot 2) = O(n) \]
\textbf{Conclusion}: We can build a heap in \textbf{linear time}.

\subsection{The Heapsort Algorithm}

Now we combine the pieces.
\begin{enumerate}
    \item Build a max-heap from the input array. ($O(n)$)
    \item The maximum element is at $A[1]$. Swap it with $A[n]$ (placing it in its final sorted position).
    \item Decrement the heap size (discarding the sorted element).
    \item The new root likely violates the heap property. Call \texttt{Max-Heapify} to fix it. ($O(\lg n)$)
    \item Repeat until the heap size is 1.
\end{enumerate}

\begin{algorithm}
    \caption{Heapsort}\label{alg:heapsort}
    \begin{algorithmic}[1]
        \Procedure{Heapsort}{$A$}
        \State \Call{Build-Max-Heap}{$A$}
        \For{$i = A.\text{length}$ \textbf{downto} $2$}
        \State Exchange $A[1]$ with $A[i]$
        \State $A.\text{heap-size} = A.\text{heap-size} - 1$
        \State \Call{Max-Heapify}{$A, 1$}
        \EndFor
        \EndProcedure
    \end{algorithmic}
\end{algorithm}

\subsubsection{Analysis}
\begin{itemize}
    \item \textbf{Build-Heap}: $O(n)$.
    \item \textbf{Loop}: Executes $n-1$ times.
    \item \textbf{Heapify}: Takes $O(\lg n)$ each time.
    \item \textbf{Total Runtime}: $O(n) + (n-1) O(\lg n) = O(n \lg n)$.
\end{itemize}

\subsubsection{Rebuild vs. Repair}
\begin{itemize}
    \item \textbf{Selection Sort} (Rebuild): After extracting the max, it sees the remaining $n-1$ elements as a chaotic bag. It spends $O(n)$ to find the next max.
    \item \textbf{Heapsort} (Repair): After extracting the max, it knows the remaining structure is \textit{almost} a valid heap. Only one path from root to leaf is violated. It spends $O(\lg n)$ to repair this specific violation.
\end{itemize}

\subsection{Priority Queues}

Heaps are the underlying data structure for efficient Priority Queues. A Max-Priority Queue supports:
\begin{itemize}
    \item \texttt{Insert(S, x)}: Insert element $x$ into set $S$. ($O(\lg n)$)
    \item \texttt{Maximum(S)}: Return the element with the largest key. ($O(1)$)
    \item \texttt{Extract-Max(S)}: Remove and return the element with the largest key. ($O(\lg n)$)
    \item \texttt{Increase-Key(S, x, k)}: Increase the value of element $x$'s key to $k$. ($O(\lg n)$ - element floats up).
\end{itemize}

\subsection{Summary}
\begin{table}[h]
    \centering
    \begin{tabular}{|l|c|c|c|c|}
        \hline
        \textbf{Algorithm} & \textbf{Time (Worst)} & \textbf{Time (Best)} & \textbf{Space} & \textbf{Stable?} \\ \hline
        Insertion Sort     & $O(n^2)$              & $O(n)$               & $O(1)$         & Yes              \\ \hline
        Merge Sort         & $O(n \lg n)$          & $O(n \lg n)$         & $O(n)$         & Yes              \\ \hline
        Heapsort           & $O(n \lg n)$          & $O(n \lg n)$         & $O(1)$         & No               \\ \hline
    \end{tabular}
    \caption{Heapsort combines the efficiency of Merge Sort with the space efficiency of Insertion Sort.}
\end{table}

\clearpage


\section{Quicksort}

\subsection{The Divide-and-Conquer Paradigm Revisited}

Quicksort is a sorting algorithm that fits within the divide-and-conquer paradigm, but it represents a philosophical mirror image to Merge Sort.

\subsubsection{Structural Comparison}
\begin{itemize}
    \item \textbf{Merge Sort}:
          \begin{itemize}
              \item \textbf{Divide}: Trivial. Simply slice the indices in half ($O(1)$).
              \item \textbf{Conquer}: Recursively sort both halves.
              \item \textbf{Combine}: \textbf{Hard}. The heavy lifting is done here, merging two sorted lists ($O(n)$).
          \end{itemize}
    \item \textbf{Quicksort}:
          \begin{itemize}
              \item \textbf{Divide}: \textbf{Hard}. The heavy lifting is done here via \texttt{Partition}. We rearrange elements relative to a pivot ($O(n)$).
              \item \textbf{Conquer}: Recursively sort the subarrays defined by the partition.
              \item \textbf{Combine}: Trivial. No work is needed; the array is sorted in place ($O(1)$).
          \end{itemize}
\end{itemize}

\subsection{Partitioning: The Core Mechanism}

The heart of Quicksort is the \texttt{Partition} procedure. It selects a \textbf{pivot} element and rearranges the array so that:
\begin{itemize}
    \item Elements smaller than or equal to the pivot are on the left.
    \item Elements larger than the pivot are on the right.
    \item The pivot is placed in its correct sorted position.
\end{itemize}

\subsubsection{The Lomuto Partition Scheme}
We maintain two pointers to define dynamic regions in the array:
\begin{itemize}
    \item \textbf{$i$}: The boundary of the "Small Region" (elements $\le x$).
    \item \textbf{$j$}: The scanning pointer (current element being inspected).
\end{itemize}

\begin{algorithm}
    \caption{Partition (A, p, r)}\label{alg:partition}
    \begin{algorithmic}[1]
        \Procedure{Partition}{$A, p, r$}
        \State $x = A[r]$ \Comment{Select the last element as pivot}
        \State $i = p - 1$ \Comment{Small region is initially empty}
        \For{$j = p$ \textbf{to} $r-1$} \Comment{Scan the array}
        \If{$A[j] \le x$}
        \State $i = i + 1$ \Comment{Expand the small region}
        \State Exchange $A[i]$ with $A[j]$ \Comment{Swap element into the small region}
        \EndIf
        \EndFor
        \State Exchange $A[i+1]$ with $A[r]$ \Comment{Place pivot between regions}
        \State \Return $i + 1$ \Comment{Return pivot index}
        \EndProcedure
    \end{algorithmic}
\end{algorithm}

\subsubsection{Trace Example}
Let's trace the execution on the array $A = [2, 8, 7, 1, 3, 5, 6, 4]$.
\begin{itemize}
    \item \textbf{Setup}: Pivot $x=4$. $i$ starts at index 0 (before array).
    \item \textbf{j=1 (2)}: $2 \le 4$. $i \to 1$. Swap $A[1] \leftrightarrow A[1]$. Array: $[2, 8, 7, \dots]$.
    \item \textbf{j=2 (8)}: $8 > 4$. No swap.
    \item \textbf{j=3 (7)}: $7 > 4$. No swap.
    \item \textbf{j=4 (1)}: $1 \le 4$. $i \to 2$. Swap $A[2] (8) \leftrightarrow A[4] (1)$. Array: $[2, 1, 7, 8, 3, \dots]$.
    \item \textbf{j=5 (3)}: $3 \le 4$. $i \to 3$. Swap $A[3] (7) \leftrightarrow A[5] (3)$. Array: $[2, 1, 3, 8, 7, \dots]$.
    \item \textbf{End}: Swap pivot $A[8] (4)$ with $A[i+1] (A[4]=8)$.
    \item \textbf{Result}: $[2, 1, 3, \mathbf{4}, 7, 5, 6, 8]$.
\end{itemize}
Observe: Left of 4 are $\{2, 1, 3\}$ (all $\le 4$). Right of 4 are $\{7, 5, 6, 8\}$ (all $> 4$).

\subsubsection{Loop Invariant Visualization}
At the start of each iteration $j$:
\begin{enumerate}
    \item $A[p \dots i] \le x$ (Known Small)
    \item $A[i+1 \dots j-1] > x$ (Known Large)
    \item $A[j \dots r-1]$ (Unknown)
    \item $A[r] = x$ (Pivot)
\end{enumerate}

\subsection{Rigorous Runtime Analysis}

\subsubsection{Worst-Case: The Unbalanced Split}
The worst case occurs when the partition routine produces one subproblem with $n-1$ elements and one with 0 elements.
\begin{itemize}
    \item \textbf{Scenario}: Input is already sorted ($1, 2, \dots, n$) or reverse sorted. Pivot (last element) is always an extreme value.
    \item \textbf{Recurrence}: $T(n) = T(n-1) + T(0) + \Theta(n) = T(n-1) + cn$.
    \item \textbf{Summation}: $T(n) = \sum_{k=1}^n ck = c \frac{n(n+1)}{2} = \Theta(n^2)$.
    \item \textbf{Visual}: The recursion tree is a "linked list" of height $n$.
\end{itemize}

\subsubsection{Best-Case: The Perfect Split}
\begin{itemize}
    \item \textbf{Scenario}: Pivot is always the median.
    \item \textbf{Recurrence}: $T(n) = 2T(n/2) + cn$.
    \item \textbf{Solution}: $\Theta(n \lg n)$ (Master Theorem Case 2).
\end{itemize}

\subsubsection{Average-Case: The Intuition of Balance}
Even a 9-to-1 split yields $O(n \lg n)$.
\[ T(n) = T(n/10) + T(9n/10) + cn \]
The tree height is $\log_{10/9} n \approx \lg n$. Since the work at each level is $\le cn$, the total is $O(n \lg n)$.

\subsubsection{Average-Case: Mathematical Proof}
Assume all permutations are equally likely. The pivot rank is uniformly distributed in $[0, n-1]$.
\[
    T(n) = \frac{1}{n} \sum_{k=0}^{n-1} (T(k) + T(n-k-1)) + \Theta(n)
\]
By symmetry:
\[
    T(n) = \frac{2}{n} \sum_{k=0}^{n-1} T(k) + \Theta(n)
\]
We suspect $T(n) \le an \lg n$. Substituting this into the sum:
\[
    \sum_{k=0}^{n-1} k \lg k \le \int_2^n x \ln x \, dx \approx \frac{n^2 \ln n}{2} - \frac{n^2}{4}
\]
Substituting back into the expression for $T(n)$ confirms the $O(n \lg n)$ bound.

\subsection{Improvements and Variants}

\subsubsection{Randomized Quicksort}
\textbf{Problem}: A fixed pivot (e.g., $A[r]$) is vulnerable to specific input patterns (like sorted arrays).
\textbf{Solution}: Pick a pivot index randomly from $[p, r]$. Swap $A[\text{random}]$ with $A[r]$ before partitioning.
\begin{itemize}
    \item \textbf{Effect}: The worst case is no longer determined by the input order, but by the random number generator (extremely unlikely).
    \item \textbf{Expected Time}: $\Theta(n \lg n)$ for \textit{any} input.
\end{itemize}

\subsubsection{Median-of-3 Partitioning}
Approximation of the true median. Select the pivot as the median of:
\[ \{ A[p], A[(p+r)/2], A[r] \} \]
This reduces the probability of bad splits and guards against sorted inputs without full randomization overhead.

\subsubsection{Handling Duplicates}
Standard Quicksort can degrade to $O(n^2)$ if all elements are equal.
\textbf{Improvement}: 3-Way Partitioning (Dutch National Flag problem).
Split array into: $\{ < x \}, \{ = x \}, \{ > x \}$.
Recursion only continues on the $< x$ and $> x$ parts.

\subsubsection{Dual-Pivot Quicksort}
(Used in Java's \texttt{Arrays.sort} for primitives).
Uses \textbf{two pivots} ($p_1, p_2$) to split the array into three segments:
\[ \{ < p_1 \}, \{ p_1 \le \dots \le p_2 \}, \{ > p_2 \} \]
This reduces the height of the recursion tree and performs fewer memory accesses on average.

\clearpage


\section{Randomisation \& Lower Bounds}

\subsection{Randomisation in Algorithm Design}

\subsubsection{The Motivation}
In previous lectures, we saw that QuickSort is efficient on average but suffers from a $\Theta(n^2)$ worst-case scenario. This worst case is triggered by specific input patterns (e.g., sorted or reverse-sorted arrays) when using a deterministic pivot (like the first or last element).

\textbf{The Problem}: A deterministic algorithm is vulnerable. A malicious adversary can construct an input that always triggers the worst case.

\textbf{The Solution}: Introduce randomness. Instead of assuming the input is random (Average Case Analysis), we make the algorithm itself random (Randomized Algorithm).
\begin{itemize}
    \item \textbf{Goal}: Ensure that the algorithm runs efficiently for \textit{all} inputs.
    \item \textbf{Philosophy}: By picking pivots randomly, we shift the dependency. The runtime no longer depends on the input order but on the sequence of random choices. No specific input can force the worst-case behavior.
\end{itemize}

\subsubsection{Randomised QuickSort}
The only change from standard QuickSort is the pivot selection. We choose a pivot uniformly at random from the subarray $A[p \dots r]$.

\begin{algorithm}
    \caption{Randomised Partition}\label{alg:randPartition}
    \begin{algorithmic}[1]
        \Procedure{Randomised-Partition}{$A, p, r$}
        \State $i = \Call{Random}{p, r}$ \Comment{Pick $i \in [p, r]$ uniformly}
        \State Exchange $A[r]$ with $A[i]$ \Comment{Move random pivot to the end}
        \State \Return \Call{Partition}{$A, p, r$} \Comment{Call standard partition}
        \EndProcedure
    \end{algorithmic}
\end{algorithm}

\begin{algorithm}
    \caption{Randomised QuickSort}\label{alg:randQuickSort}
    \begin{algorithmic}[1]
        \Procedure{Randomised-QuickSort}{$A, p, r$}
        \If{$p < r$}
        \State $q = \Call{Randomised-Partition}{A, p, r}$
        \State \Call{Randomised-QuickSort}{$A, p, q-1$}
        \State \Call{Randomised-QuickSort}{$A, q+1, r$}
        \EndIf
        \EndProcedure
    \end{algorithmic}
\end{algorithm}

\subsection{Analysis of Expected Runtime}

For a randomized algorithm, we analyze the \textbf{Expected Running Time} $E[T(n)]$. We count the number of comparisons, as this dominates the runtime.

\subsubsection{Setup: Indicator Random Variables}
Let $X$ be the total number of comparisons.
Let the sorted elements of the array be $z_1 < z_2 < \dots < z_n$.
We define an indicator random variable $X_{ij}$ for any pair $i < j$:
\[
    X_{ij} = \mathbb{I}\{ z_i \text{ is compared to } z_j \} = \begin{cases}
        1 & \text{if } z_i \text{ and } z_j \text{ are compared} \\
        0 & \text{otherwise}
    \end{cases}
\]
The total number of comparisons is the sum of all pairs:
\[ X = \sum_{i=1}^{n-1} \sum_{j=i+1}^{n} X_{ij} \]

\subsubsection{Linearity of Expectation}
A powerful property of expectation is that the expectation of a sum is the sum of expectations, \textit{even if the variables are dependent}.
\[ E[X] = E\left[ \sum_{i=1}^{n-1} \sum_{j=i+1}^{n} X_{ij} \right] = \sum_{i=1}^{n-1} \sum_{j=i+1}^{n} E[X_{ij}] \]
Since $E[X_{ij}] = \Pr(z_i \text{ is compared to } z_j)$, we just need to find this probability.

\subsubsection{The Probability of Comparison}
\textbf{Critical Insight}: When are two elements $z_i$ and $z_j$ compared?
\begin{itemize}
    \item Elements are compared \textbf{only} when one of them is chosen as a pivot.
    \item Once a pivot $x$ is chosen, all other elements are compared to $x$ and then separated into different sets ($<x$ and $>x$).
    \item Consider the set of elements $Z_{ij} = \{z_i, z_{i+1}, \dots, z_j\}$.
    \item If any element $x \in Z_{ij}$ such that $z_i < x < z_j$ is chosen as a pivot \textit{before} $z_i$ or $z_j$, then $z_i$ and $z_j$ will be separated into different subproblems and \textbf{never compared}.
    \item Thus, $z_i$ and $z_j$ are compared if and only if the \textit{first} element selected as a pivot from the set $Z_{ij}$ is either $z_i$ or $z_j$.
\end{itemize}

The set $Z_{ij}$ has $j - i + 1$ elements. Since pivots are chosen uniformly at random, every element in $Z_{ij}$ has an equal chance ($1 / (j-i+1)$) of being the first one picked.
There are 2 favorable outcomes ($z_i$ or $z_j$).
\[ \Pr(z_i \text{ is compared to } z_j) = \frac{2}{j - i + 1} \]

\subsubsection{Summation and Result}
Substituting this back into the expectation sum:
\[ E[X] = \sum_{i=1}^{n-1} \sum_{j=i+1}^{n} \frac{2}{j - i + 1} \]
Let $k = j - i$ (the distance between elements). As $j$ goes from $i+1$ to $n$, $k$ goes from $1$ to $n-i$.
\[ E[X] = \sum_{i=1}^{n-1} \sum_{k=1}^{n-i} \frac{2}{k+1} < \sum_{i=1}^{n-1} \sum_{k=1}^{n} \frac{2}{k} \]
The inner sum is the Harmonic series $H_n \approx \ln n$.
\[ E[X] < \sum_{i=1}^{n-1} 2 \ln n = 2(n-1) \ln n = O(n \log n) \]
\textbf{Conclusion}: The expected runtime of Randomized QuickSort is $O(n \log n)$.

\subsection{Lower Bounds for Comparison Sorts}

We have seen several algorithms (MergeSort, HeapSort) with $O(n \log n)$ worst-case time. Can we do better? Can we sort in $O(n)$?

\subsubsection{The Comparison Model}
We restrict our analysis to \textbf{Comparison Sorts}: algorithms that only gain information about the input sequence by comparing pairs of elements ($a_i < a_j$, etc.). We cannot inspect the values themselves or use them as array indices.

\subsubsection{Decision Trees}
Any comparison sort can be modeled abstractly as a \textbf{Decision Tree}.
\begin{itemize}
    \item \textbf{Root}: The first comparison made by the algorithm.
    \item \textbf{Internal Nodes}: Comparisons $a_i : a_j$.
    \item \textbf{Edges}: The result of the comparison (Left: $\le$, Right: $>$).
    \item \textbf{Leaves}: A permutation of the input elements (the final sorted order).
\end{itemize}
An execution of the algorithm corresponds to a path from the root to a leaf. The length of the path is the number of comparisons. The \textbf{height} of the tree is the worst-case number of comparisons.



\subsubsection{The Lower Bound Theorem}
\begin{theorem}
    Any comparison sort algorithm requires $\Omega(n \log n)$ comparisons in the worst case.
\end{theorem}
\begin{proof}
    Let the decision tree have height $h$ and $L$ leaves.
    \begin{enumerate}
        \item \textbf{Necessary Leaves}: There are $n!$ possible permutations of $n$ distinct elements. For the algorithm to be correct, it must be able to produce \textit{any} of these $n!$ permutations as output. Thus, the tree must have at least $n!$ reachable leaves:
              \[ L \ge n! \]
        \item \textbf{Binary Tree Property}: A binary tree of height $h$ has at most $2^h$ leaves.
              \[ L \le 2^h \]
        \item \textbf{Combining}:
              \[ 2^h \ge n! \]
              Taking logarithms base 2:
              \[ h \ge \lg(n!) \]
        \item \textbf{Stirling's Approximation}: For large $n$, $n! \approx \sqrt{2\pi n} (\frac{n}{e})^n$.
              \[ \lg(n!) \approx \lg\left(\left(\frac{n}{e}\right)^n\right) = n \lg n - n \lg e = \Theta(n \log n) \]
        \item \textbf{Conclusion}: $h = \Omega(n \log n)$.
    \end{enumerate}
\end{proof}

\subsection{Summary of Sorting Algorithms}

We can now classify sorting algorithms based on this lower bound.

\begin{table}[h]
    \centering
    \begin{tabular}{|l|c|c|c|c|}
        \hline
        \textbf{Algorithm} & \textbf{Best Time} & \textbf{Avg Time} & \textbf{Worst Time} & \textbf{Space} \\ \hline
        Insertion Sort     & $O(n)$             & $O(n^2)$          & $O(n^2)$            & $O(1)$         \\ \hline
        Merge Sort         & $O(n \lg n)$       & $O(n \lg n)$      & $O(n \lg n)$        & $O(n)$         \\ \hline
        Heap Sort          & $O(n \lg n)$       & $O(n \lg n)$      & $O(n \lg n)$        & $O(1)$         \\ \hline
        Quick Sort         & $O(n \lg n)$       & $O(n \lg n)$      & $O(n^2)$            & $O(\lg n)$     \\ \hline
        Rand. Quick Sort   & $O(n \lg n)$       & $O(n \lg n)$      & $O(n^2)$            & $O(\lg n)$     \\ \hline
    \end{tabular}
    \caption{Comparison of Sorting Algorithms. Note: Heap Sort and Merge Sort are asymptotically optimal comparison sorts.}
\end{table}

\textbf{Note on Linear Time}: The $\Omega(n \log n)$ bound applies only to \textit{comparison} sorts. If we know more about the input (e.g., integers in a small range), we can break this barrier using algorithms like Counting Sort or Radix Sort (next section).

\clearpage


\section{Sorting in Linear Time}

\subsection{Breaking the Speed Limit}

\subsubsection{The Comparison Bottleneck}
In previous lectures, we established a fundamental limit: \textbf{Any comparison-based sorting algorithm requires $\Omega(n \log n)$ time in the worst case.} This is a mathematical certainty derived from the height of the decision tree.

\subsubsection{The Way Out}
How can we sort faster than $O(n \log n)$? We must change the rules of the game.
\begin{itemize}
    \item \textbf{Stop Comparing}: Do not ask "Is $A[i] < A[j]$?".
    \item \textbf{Start Inspecting}: Use the actual numerical value of the keys to determine their position directly.
    \item \textbf{Constraint}: These algorithms assume the input comes from a restricted domain (e.g., small integers).
\end{itemize}

\subsection{Counting Sort}

Counting Sort is the foundational algorithm for linear-time sorting. It relies on the assumption that input elements are integers in the range $\{0, \dots, k\}$.

\subsubsection{The Algorithm logic}
The core idea is determining, for each input element $x$, exactly how many elements are less than or equal to $x$. If there are 17 elements smaller than $x$, then $x$ belongs at position 18.

\begin{algorithm}
    \caption{Counting Sort}\label{alg:countingSort}
    \begin{algorithmic}[1]
        \Procedure{Counting-Sort}{$A, B, k$}
        \State Let $C[0 \dots k]$ be a new array initialized to 0
        \State \Comment{Step 1: Frequency Count (Histogram)}
        \For{$j = 1$ \textbf{to} $A.\text{length}$}
        \State $C[A[j]] = C[A[j]] + 1$
        \EndFor
        \State \Comment{Step 2: Cumulative Sums (Prefix Sums)}
        \For{$i = 1$ \textbf{to} $k$}
        \State $C[i] = C[i] + C[i-1]$
        \EndFor
        \State \Comment{$C[i]$ now contains the number of elements $\le i$}
        \State \Comment{Step 3: Placement (Reverse Traversal)}
        \For{$j = A.\text{length}$ \textbf{downto} 1}
        \State $B[C[A[j]]] = A[j]$ \Comment{Place element at its correct sorted position}
        \State $C[A[j]] = C[A[j]] - 1$ \Comment{Decrement count for the next instance}
        \EndFor
        \EndProcedure
    \end{algorithmic}
\end{algorithm}

\subsubsection{Why Traverse Backwards? (Stability)}
Notice the loop in Step 3 goes \textbf{downto} 1. This is crucial for \textbf{Stability}.
\begin{itemize}
    \item Suppose the input $A$ contains two copies of value 3: $3_a$ at index 2 and $3_b$ at index 5.
    \item The algorithm encounters $3_b$ first. It places $3_b$ at position $C[3]$ and decrements $C[3]$.
    \item Later, it encounters $3_a$. It places $3_a$ at position $C[3]$ (which is now one less).
    \item Result: $3_b$ appears after $3_a$ in the output, preserving their original relative order.
\end{itemize}

\subsubsection{Complexity Analysis}
\begin{itemize}
    \item \textbf{Time}: $\Theta(k)$ (init) + $\Theta(n)$ (count) + $\Theta(k)$ (sum) + $\Theta(n)$ (place) = $\Theta(n + k)$.
    \item \textbf{Space}: $\Theta(n + k)$ for arrays $B$ and $C$.
    \item \textbf{Implication}: If $k = O(n)$ (the range is linear in the number of elements), Counting Sort runs in \textbf{linear time} $\Theta(n)$.
\end{itemize}

\subsection{Radix Sort}

Counting Sort is fast but impractical for large ranges (e.g., sorting 32-bit integers would require $k=2^{32}$ memory). Radix Sort solves this by sorting digit-by-digit.

\subsubsection{The Algorithm (LSD Approach)}
We sort by the Least Significant Digit (LSD) first, then the next digit, and so on, up to the Most Significant Digit (MSD).
\textbf{Crucial Requirement}: The sort used for each digit must be \textbf{stable}.

\begin{algorithm}
    \caption{Radix Sort}\label{alg:radixSort}
    \begin{algorithmic}[1]
        \Procedure{Radix-Sort}{$A, d$}
        \For{$i = 1$ \textbf{to} $d$}
        \State Use a \textbf{stable sort} to sort array $A$ on digit $i$
        \EndFor
        \EndProcedure
    \end{algorithmic}
\end{algorithm}

\subsubsection{Why LSD works (Intuition)}
Why start from the right?
\begin{itemize}
    \item After sorting on digit 1, the array is sorted by the last digit.
    \item When we sort on digit 2, stability ensures that if two numbers have the same 2nd digit, their order determined by the 1st digit is preserved.
    \item \textbf{Inductive Hypothesis}: After sorting on digit $i$, the array is sorted according to the value of the number formed by the last $i$ digits.
\end{itemize}

\subsubsection{Trace Example}
Sorting: $[329, 457, 657, 839, 436, 720, 355]$
\begin{enumerate}
    \item \textbf{Sort on ones digit}:
          $72\mathbf{0}, 35\mathbf{5}, 43\mathbf{6}, 45\mathbf{7}, 65\mathbf{7}, 32\mathbf{9}, 83\mathbf{9}$
          (Note: 457 comes before 657 because it appeared first; 329 before 839).
    \item \textbf{Sort on tens digit}:
          $7\mathbf{2}0, 3\mathbf{2}9, 4\mathbf{3}6, 8\mathbf{3}9, 3\mathbf{5}5, 4\mathbf{5}7, 6\mathbf{5}7$
    \item \textbf{Sort on hundreds digit}:
          $\mathbf{3}29, \mathbf{3}55, \mathbf{4}36, \mathbf{4}57, \mathbf{6}57, \mathbf{7}20, \mathbf{8}39$
          Sorted!
\end{enumerate}

\subsection{Advanced Analysis: Breaking the $O(n^3)$ Range}

Consider the problem: Sort $n$ integers in the range $0$ to $n^3 - 1$.

\subsubsection{Attempt 1: Comparison Sort}
Standard Merge/Quick Sort takes $\Theta(n \log n)$. Good, but not linear.

\subsubsection{Attempt 2: Counting Sort}
The range is $k = n^3$.
Runtime: $\Theta(n + k) = \Theta(n + n^3) = \Theta(n^3)$. This is terrible.

\subsubsection{Attempt 3: Radix Sort (Base 10)}
Range $M = n^3$. Number of digits $d = \log_{10} M = 3 \log_{10} n$.
Runtime: $\Theta(d \cdot n) = \Theta(n \log n)$. Still not linear.

\subsubsection{Attempt 4: Radix Sort (Base $n$)}
Treat the numbers as being written in \textbf{base $n$}.
\begin{itemize}
    \item The "digits" now range from $0$ to $n-1$. So $k = n$.
    \item How many digits $d$?
          \[ d = \log_n (n^3 - 1) \approx \log_n (n^3) = 3 \]
          So we only have 3 "digits".
    \item \textbf{Runtime}:
          \[ T(n) = \Theta(d(n + k)) = \Theta(3(n + n)) = \Theta(6n) = \Theta(n) \]
\end{itemize}
\textbf{Conclusion}: By changing the base to $n$, we can sort integers up to $n^k$ (for any constant $k$) in linear time.

\subsection{Summary Table}

\begin{table}[h]
    \centering
    \begin{tabular}{|l|c|c|l|}
        \hline
        \textbf{Algorithm} & \textbf{Time}      & \textbf{Space} & \textbf{Constraints}                   \\ \hline
        Counting Sort      & $\Theta(n+k)$      & $\Theta(n+k)$  & Efficient when $k = O(n)$.             \\ \hline
        Radix Sort         & $\Theta(d(n+k))$   & $\Theta(n+k)$  & Efficient when numbers have fixed $d$. \\ \hline
        Comparison Sorts   & $\Omega(n \log n)$ & Varies         & General purpose.                       \\ \hline
    \end{tabular}
    \caption{Comparison of Linear-Time vs General Sorts}
\end{table}

\clearpage


\section{Elementary Data Structures}

\subsection{Foundations: The Limitations of Arrays}

\subsubsection{Data Structures vs. Raw Memory}
A \textbf{Data Structure} is not merely a container; it is a strategic way of organizing a \textbf{finite dynamic set} of elements to optimize specific operations.
Elements usually consist of:
\begin{itemize}
    \item \textbf{Key}: The identifier (used for sorting/searching).
    \item \textbf{Satellite Data}: The payload.
    \item \textbf{Attributes}: Meta-data maintained by the structure (e.g., \texttt{next}, \texttt{top}, \texttt{head}).
\end{itemize}

\subsubsection{The Array Bottleneck}
Why do we need anything other than arrays? Arrays map indices to memory addresses directly ($O(1)$ access), but they are \textbf{rigid}.
\begin{itemize}
    \item \textbf{Unsorted Array}:
          \begin{itemize}
              \item \textbf{Insertion}: $O(1)$ (append to end).
              \item \textbf{Search}: $\Theta(n)$ (linear scan).
          \end{itemize}
    \item \textbf{Sorted Array}:
          \begin{itemize}
              \item \textbf{Search}: $O(\log n)$ (Binary Search).
              \item \textbf{Insertion/Deletion}: $\Theta(n)$. \textbf{Why?} Because elements are stored contiguously in memory. Inserting $x$ at index $i$ requires shifting elements $i \dots n$ one position to the right.
          \end{itemize}
\end{itemize}
\textbf{Conclusion}: Arrays cannot simultaneously offer fast search \textbf{and} fast dynamic updates.

\subsection{Stacks (LIFO)}

\subsubsection{Philosophy}
A \textbf{Stack} enforces a \textbf{Last-In, First-Out (LIFO)} policy. By restricting access \textbf{strictly} to the "top" element, we eliminate the need for shifting. Thus, operations become $O(1)$.

\subsubsection{Array Implementation}
We can simulate a stack using an array $S$ and a pointer $S.top$.
\begin{itemize}
    \item \textbf{Empty}: $S.top = 0$.
    \item \textbf{Push}: Increment $top$, then write. (Check for Overflow).
    \item \textbf{Pop}: Read, then decrement $top$. (Check for Underflow).
\end{itemize}

\begin{algorithm}
    \caption{Stack Operations}\label{alg:stack}
    \begin{algorithmic}[1]
        \Procedure{Push}{$S, x$}
        \If{$S.top == S.size$}
        \State \textbf{Error} "Overflow"
        \Else
        \State $S.top = S.top + 1$
        \State $S[S.top] = x$
        \EndIf
        \EndProcedure
        \Statex
        \Procedure{Pop}{$S$}
        \If{\Call{Stack-Empty}{$S$}}
        \State \textbf{Error} "Underflow"
        \Else
        \State $S.top = S.top - 1$
        \State \textbf{Return} $S[S.top + 1]$
        \EndIf
        \EndProcedure
    \end{algorithmic}
\end{algorithm}

\subsubsection{Algorithmic Applications}
Stacks are the natural data structure for processing \textbf{nested} or \textbf{recursive} structures.

\textbf{1. Bracket Balance Checking}:
Problem: Is `{[()][()]}` valid?
\begin{itemize}
    \item \textbf{Algorithm}: Iterate through the string.
    \item If \textbf{Opening} (`(`, `[`, `{`): \textbf{Push} onto stack.
    \item If \textbf{Closing} (`)`, `]`, `}`): \textbf{Pop}. Check if the popped element matches the current closing bracket.
    \item \textbf{Result}: Valid iff stack is empty at the end and no mismatches occurred.
\end{itemize}

\textbf{2. Postfix Expression Evaluation}:
Infix: $(9+3) * (4*2)$ $\rightarrow$ Postfix: $9 \ 3 \ + \ 4 \ 2 \ * \ *$
\begin{itemize}
    \item \textbf{Algorithm}:
          \begin{itemize}
              \item Operand ($9, 3$): \textbf{Push}.
              \item Operator ($+$): \textbf{Pop} $b$ ($3$), \textbf{Pop} $a$ ($9$). Compute $a+b$. \textbf{Push} result ($12$).
          \end{itemize}
    \item \textbf{Advantage}: Eliminates need for parentheses and operator precedence rules.
\end{itemize}

\subsection{Queues (FIFO)}

\subsubsection{Philosophy}
A \textbf{Queue} enforces a \textbf{First-In, First-Out (FIFO)} policy. It models fairness (lines, buffers).
\begin{itemize}
    \item \textbf{Enqueue}: Add to \textbf{Tail}.
    \item \textbf{Dequeue}: Remove from \textbf{Head}.
\end{itemize}

\subsubsection{Circular Buffer Implementation}
To implement a queue in a fixed array without shifting elements after every dequeue, we use \textbf{modular arithmetic} to wrap indices around.
\begin{itemize}
    \item $Q.head$: Index of the element to dequeue.
    \item $Q.tail$: Index of the \textbf{next empty slot}.
    \item \textbf{Full Condition}: $(Q.tail + 1) == Q.head$ (conceptually).
\end{itemize}

\begin{algorithm}
    \caption{Queue Operations}\label{alg:queue}
    \begin{algorithmic}[1]
        \Procedure{Enqueue}{$Q, x$}
        \State $Q[Q.tail] = x$
        \If{$Q.tail == Q.size$} $Q.tail = 1$ \Comment{Wrap around}
        \Else \ $Q.tail = Q.tail + 1$
        \EndIf
        \EndProcedure
        \Statex
        \Procedure{Dequeue}{$Q$}
        \State $x = Q[Q.head]$
        \If{$Q.head == Q.size$} $Q.head = 1$ \Comment{Wrap around}
        \Else \ $Q.head = Q.head + 1$
        \EndIf
        \State \textbf{Return} $x$
        \EndProcedure
    \end{algorithmic}
\end{algorithm}
\textbf{Runtime}: $O(1)$ for both.

\subsection{Priority Queues}

Stacks and Queues are determined by \textit{time}. \textbf{Priority Queues} are determined by \textit{importance} (Key).
\begin{itemize}
    \item \textbf{Operations}: \texttt{Insert}, \texttt{Extract-Max}.
    \item \textbf{Implementation}:
          \begin{itemize}
              \item Sorted Array: Extract-Max is $O(1)$, but Insert is $O(n)$.
              \item \textbf{Heap}: Both operations are $O(\log n)$. (Covered in detail in Heapsort lecture).
          \end{itemize}
\end{itemize}

\subsection{Linked Lists}

\subsubsection{Breaking the Contiguity Constraint}
Linked Lists decouple logical order from physical memory order.
\begin{itemize}
    \item \textbf{Pros}: Dynamic size; $O(1)$ insertion/deletion (if location is known).
    \item \textbf{Cons}: No random access; $\Theta(n)$ search; Extra memory for pointers.
\end{itemize}

\subsubsection{Operations Analysis}

\textbf{1. Searching}:
Must scan sequentially from \texttt{head}.
\[ T(n) = \Theta(n) \]

\textbf{2. Insertion (at front)}:
\begin{itemize}
    \item New node points to old head.
    \item Old head's \texttt{prev} points to new node.
    \item Update \texttt{head} pointer.
\end{itemize}
\[ T(n) = O(1) \]

\textbf{3. Deletion}:
This operation highlights the difference between "Knowing the Key" and "Knowing the Node".
\begin{itemize}
    \item \textbf{Case A: Given Key $k$}:
          We must first \texttt{Search} for $k$ ($\Theta(n)$), then delete. Total: $\Theta(n)$.
    \item \textbf{Case B: Given Pointer $x$}:
          We simply rewire the pointers of $x.prev$ and $x.next$ to bypass $x$.
          \[ x.prev.next = x.next \]
          \[ x.next.prev = x.prev \]
          Total: $O(1)$.
\end{itemize}

\subsection{Summary of Complexities}

\begin{table}[h]
    \centering
    \begin{tabular}{|l|c|c|c|l|}
        \hline
        \textbf{Data Structure} & \textbf{Insert} & \textbf{Delete} & \textbf{Search} & \textbf{Constraint} \\ \hline
        Array (Unsorted)        & $O(1)$          & $O(1)^*$        & $\Theta(n)$     & Static Size         \\ \hline
        Array (Sorted)          & $\Theta(n)$     & $\Theta(n)$     & $O(\log n)$     & Static Size         \\ \hline
        Stack                   & $O(1)$          & $O(1)$          & -               & LIFO Only           \\ \hline
        Queue                   & $O(1)$          & $O(1)$          & -               & FIFO Only           \\ \hline
        Linked List             & $O(1)$          & $O(1)^{**}$     & $\Theta(n)$     & Sequential Access   \\ \hline
    \end{tabular}
    \caption{*Delete by swapping with last element. **Delete given pointer to node.}
\end{table}

\clearpage


\section{Binary Search Trees}

\subsection{Introduction to Trees}

\subsubsection{Definitions and Terminology}

\begin{definition}[Binary Tree]\index{Binary Tree}
    A \textbf{binary tree} is a structure defined recursively. It is a finite set of nodes that is either:
    \begin{itemize}
        \item Empty (contains no nodes).
        \item Consists of a \textbf{root} node and two disjoint binary trees called the \textbf{left subtree} and the \textbf{right subtree}.
    \end{itemize}
\end{definition}

Key terminology:
\begin{itemize}
    \item \textbf{Leaf}: A node with no children (left and right subtrees are empty).
    \item \textbf{Height of a node}: The length (number of edges) of the longest simple path from the node to a leaf.
    \item \textbf{Height of a tree}: The height of the root.
    \item \textbf{Depth of a node}: The length of the path from the root to the node.
    \item \textbf{Full Binary Tree}: A binary tree in which every node is either a leaf or has exactly two children.
\end{itemize}

\subsubsection{Inductive Proofs on Trees}
Because trees are defined recursively, they are well-suited for proofs by induction.

\begin{theorem}
    A binary tree of height $h$ has at most $2^h$ leaves.
\end{theorem}
\begin{proof}
    We proceed by induction on the height $h$.
    \begin{itemize}
        \item \textbf{Base Case ($h=0$)}: The tree consists of a single root node. It has $1$ leaf. Since $2^0 = 1$, the base case holds.
        \item \textbf{Inductive Hypothesis}: Assume that for any binary tree of height $h-1$, the number of leaves is at most $2^{h-1}$.
        \item \textbf{Inductive Step}: Consider a tree $T$ of height $h$. Its root has two subtrees (left and right), each with height at most $h-1$.
              Let $L(T)$ be the number of leaves. $L(T) = L(T_{left}) + L(T_{right})$.
              By the hypothesis, $L(T_{left}) \le 2^{h-1}$ and $L(T_{right}) \le 2^{h-1}$.
              Thus, $L(T) \le 2^{h-1} + 2^{h-1} = 2 \cdot 2^{h-1} = 2^h$.
    \end{itemize}
\end{proof}

\subsection{Binary Search Trees (BST)}

\begin{definition}[Binary Search Tree Property]\index{BST Property}
    Let $x$ be a node in a binary search tree.
    \begin{itemize}
        \item If $y$ is a node in the \textbf{left} subtree of $x$, then $y.key \le x.key$.
        \item If $y$ is a node in the \textbf{right} subtree of $x$, then $y.key \ge x.key$.
    \end{itemize}
\end{definition}

This property allows us to print all keys in sorted order using a simple recursive algorithm.

\subsubsection{Tree Walks}

\begin{algorithm}
    \caption{Inorder Tree Walk}\label{alg:inorder}
    \begin{algorithmic}[1]
        \Procedure{Inorder-Walk}{$x$}
        \If{$x \ne \text{NIL}$}
        \State \Call{Inorder-Walk}{$x.left$}
        \State Print $x.key$
        \State \Call{Inorder-Walk}{$x.right$}
        \EndIf
        \EndProcedure
    \end{algorithmic}
\end{algorithm}

\begin{theorem}
    If $x$ is the root of an $n$-node subtree, the call \texttt{Inorder-Walk(x)} takes $\Theta(n)$ time.
\end{theorem}
\begin{proof}
    We can prove this using the \textbf{Accounting Method}. Assign a constant cost $c$ to each node. The procedure visits each node exactly once (when printing), and traverses each edge exactly twice (once going down, once returning up). Since the number of edges is $n-1$, the total work is proportional to the number of nodes. Thus, $T(n) = \Theta(n)$.
\end{proof}

\subsection{Query Operations}

Operations on a BST generally depend on the height $h$ of the tree, taking $O(h)$ time.

\subsubsection{Search}
To find a node with key $k$, we trace a path from the root.
\begin{algorithm}
    \caption{Tree Search}\label{alg:treeSearch}
    \begin{algorithmic}[1]
        \Procedure{Tree-Search}{$x, k$}
        \If{$x == \text{NIL}$ \textbf{or} $k == x.key$}
        \State \Return $x$
        \EndIf
        \If{$k < x.key$}
        \State \Return \Call{Tree-Search}{$x.left, k$}
        \Else
        \State \Return \Call{Tree-Search}{$x.right, k$}
        \EndIf
        \EndProcedure
    \end{algorithmic}
\end{algorithm}

\subsubsection{Minimum and Maximum}
Due to the BST property:
\begin{itemize}
    \item The \textbf{minimum} element is found by following \texttt{left} pointers until a node with no left child is reached.
    \item The \textbf{maximum} element is found by following \texttt{right} pointers until a node with no right child is reached.
\end{itemize}
\textbf{Runtime}: $O(h)$.

\subsubsection{Successor}
The \textbf{successor} of a node $x$ is the node with the smallest key greater than $x.key$.
\begin{itemize}
    \item \textbf{Case 1}: If the right subtree of $x$ is non-empty, the successor is the \textbf{minimum} node in the right subtree.
    \item \textbf{Case 2}: If the right subtree is empty, the successor is the lowest ancestor of $x$ whose left child is also an ancestor of $x$ (i.e., we go up until we turn right).
\end{itemize}

\begin{algorithm}
    \caption{Tree Successor}\label{alg:successor}
    \begin{algorithmic}[1]
        \Procedure{Tree-Successor}{$x$}
        \If{$x.right \ne \text{NIL}$}
        \State \Return \Call{Tree-Minimum}{$x.right$}
        \EndIf
        \State $y = x.p$
        \While{$y \ne \text{NIL}$ \textbf{and} $x == y.right$} \Comment{Go up while we are a right child}
        \State $x = y$
        \State $y = y.p$
        \EndWhile
        \State \Return $y$
        \EndProcedure
    \end{algorithmic}
\end{algorithm}

\subsection{Modifying Operations}

\subsubsection{Insertion}
To insert a new value $z$, we trace a path from the root downwards, similar to \texttt{Tree-Search}. We maintain a "trailing pointer" $y$ to keep track of the parent. When we hit \texttt{NIL}, we attach $z$ to $y$.

\begin{algorithm}
    \caption{Tree Insert}\label{alg:treeInsert}
    \begin{algorithmic}[1]
        \Procedure{Tree-Insert}{$T, z$}
        \State $y = \text{NIL}$
        \State $x = T.root$
        \While{$x \ne \text{NIL}$}
        \State $y = x$
        \If{$z.key < x.key$} $x = x.left$
        \Else \ $x = x.right$
        \EndIf
        \EndWhile
        \State $z.p = y$
        \If{$y == \text{NIL}$} $T.root = z$ \Comment{Tree was empty}
        \ElsIf{$z.key < y.key$} $y.left = z$
        \Else \ $y.right = z$
        \EndIf
        \EndProcedure
    \end{algorithmic}
\end{algorithm}
\textbf{Runtime}: $O(h)$.

\subsubsection{Deletion}
Deleting a node $z$ is the most complex operation. We handle three cases:
\begin{enumerate}
    \item \textbf{No children}: $z$ is a leaf. Simply remove it.
    \item \textbf{One child}: Splice $z$ out. Replace $z$ with its child.
    \item \textbf{Two children}: Find $z$'s successor $y$. $y$ must lie in $z$'s right subtree and has no left child. We replace $z$'s content with $y$'s content and splice $y$ out.
\end{enumerate}

We use a subroutine \texttt{Transplant} to replace one subtree rooted at $u$ with another rooted at $v$.

\begin{algorithm}
    \caption{Transplant}\label{alg:transplant}
    \begin{algorithmic}[1]
        \Procedure{Transplant}{$T, u, v$}
        \If{$u.p == \text{NIL}$}
        \State $T.root = v$
        \ElsIf{$u == u.p.left$}
        \State $u.p.left = v$
        \Else
        \State $u.p.right = v$
        \EndIf
        \If{$v \ne \text{NIL}$}
        \State $v.p = u.p$
        \EndIf
        \EndProcedure
    \end{algorithmic}
\end{algorithm}

\begin{algorithm}
    \caption{Tree Delete}\label{alg:treeDelete}
    \begin{algorithmic}[1]
        \Procedure{Tree-Delete}{$T, z$}
        \If{$z.left == \text{NIL}$} \Comment{Case 1 or Case 2 (no left child)}
        \State \Call{Transplant}{$T, z, z.right$}
        \ElsIf{$z.right == \text{NIL}$} \Comment{Case 2 (no right child)}
        \State \Call{Transplant}{$T, z, z.left$}
        \Else \Comment{Case 3: Two children}
        \State $y = \Call{Tree-Minimum}{z.right}$ \Comment{Find successor}
        \If{$y.p \ne z$} \Comment{If successor is not immediate child}
        \State \Call{Transplant}{$T, y, y.right$} \Comment{Replace y with its right child}
        \State $y.right = z.right$
        \State $y.right.p = y$
        \EndIf
        \State \Call{Transplant}{$T, z, y$} \Comment{Replace z with y}
        \State $y.left = z.left$
        \State $y.left.p = y$
        \EndIf
        \EndProcedure
    \end{algorithmic}
\end{algorithm}
\textbf{Runtime}: $O(h)$ because finding the successor and handling pointers takes time proportional to height.

\subsection{Performance Analysis}

The runtime of BST operations is determined by the height $h$.
\begin{itemize}
    \item \textbf{Worst Case}: If we insert elements in sorted (or reverse sorted) order, the tree becomes a linear chain.
          \[ h = n \implies \text{Runtime} = \Theta(n) \]
    \item \textbf{Best/Average Case}: In a randomly built BST, the expected height is logarithmic.
          \[ h = O(\lg n) \implies \text{Runtime} = O(\lg n) \]
\end{itemize}

This instability motivates the need for \textbf{Balanced Binary Search Trees} (like AVL trees or Red-Black trees), which guarantee $h = O(\lg n)$.

\clearpage


\section{AVL Trees}

\subsection{Motivation}

Standard Binary Search Trees (BSTs) support operations (Search, Insert, Delete, Min, Max) in $O(h)$ time. However, in the worst case (e.g., inserting sorted data), the height $h$ can degrade to $O(n)$, making operations linear.
To guarantee $O(\lg n)$ performance, we need to enforce a balance condition. \textbf{AVL Trees} (Adelson-Velsky and Landis, 1962) are the first self-balancing BSTs invented.

\subsection{Definition and Properties}

\subsubsection{AVL Property}

\begin{definition}[AVL Tree]\index{AVL Tree}
    An \textbf{AVL tree} is a binary search tree that satisfies the \textbf{AVL Property}:
    For every node $v$ in the tree, the heights of its left and right subtrees differ by at most 1.
\end{definition}

\begin{definition}[Balance Factor]\index{Balance Factor}
    The \textbf{balance factor} of a node $v$, denoted $bal(v)$, is defined as the height of the left subtree minus the height of the right subtree:
    \[ bal(v) = height(v.left) - height(v.right) \]
    In an AVL tree, for every node $v$, $bal(v) \in \{-1, 0, 1\}$.
\end{definition}

\subsubsection{Height Analysis}

\begin{theorem}
    The height $h$ of an AVL tree with $n$ nodes is $O(\lg n)$. More specifically, $h < 1.44 \lg (n+2)$.
\end{theorem}
\begin{proof}
    Let $N(h)$ be the minimum number of nodes in an AVL tree of height $h$.
    \begin{itemize}
        \item $N(0) = 1$ (Root only, height 0).
        \item $N(1) = 2$ (Root + 1 child).
        \item For general $h$, the root must have two subtrees. To minimize nodes while maintaining height $h$, one subtree must have height $h-1$ and the other $h-2$.
              \[ N(h) = 1 + N(h-1) + N(h-2) \]
    \end{itemize}
    This recurrence is strictly greater than the Fibonacci sequence ($F_h \approx \phi^h / \sqrt{5}$ where $\phi \approx 1.618$).
    Solving the recurrence implies $n > \phi^h$, thus $h < \log_\phi n \approx 1.44 \lg n$.
\end{proof}

\subsection{Rotations}

Rotations are the fundamental operations to rebalance a tree without violating the BST property (inorder traversal remains unchanged). They take $O(1)$ time.

\subsubsection{Right Rotation}
Used when the left subtree is too heavy.
\begin{algorithm}
    \caption{Right Rotation}\label{alg:rightRotate}
    \begin{algorithmic}[1]
        \Procedure{Right-Rotate}{$T, y$}
        \State $x = y.left$
        \State $y.left = x.right$ \Comment{Turn x's right subtree into y's left subtree}
        \If{$x.right \ne \text{NIL}$} $x.right.p = y$
        \EndIf
        \State $x.p = y.p$ \Comment{Link x to y's parent}
        \If{$y.p == \text{NIL}$} $T.root = x$
        \ElsIf{$y == y.p.right$} $y.p.right = x$
        \Else \ $y.p.left = x$
        \EndIf
        \State $x.right = y$ \Comment{Put y on x's right}
        \State $y.p = x$
        \State Update heights of $x$ and $y$
        \EndProcedure
    \end{algorithmic}
\end{algorithm}

\subsubsection{Left Rotation}
Used when the right subtree is too heavy. Symmetric to Right Rotation.



\subsection{Insertion}

\textbf{Goal}: Insert a node $z$ and restore the AVL property if violated.
\textbf{Steps}:
\begin{enumerate}
    \item Perform a standard BST insertion.
    \item Update the height of ancestors of $z$.
    \item Check for imbalance: If a node $v$ has $|bal(v)| > 1$, perform rotations.
\end{enumerate}

Let $v$ be the lowest unbalanced node (the first one encountered moving up from $z$). Let $x$ be the child of $v$ in the direction of the imbalance. There are 4 cases:

\subsubsection{Case 1: Left-Left (LL)}
\begin{itemize}
    \item \textbf{Condition}: $v$ is heavy on the left ($bal(v) = +2$) AND $x$ is heavy on the left ($bal(x) = +1$ or $0$).
    \item \textbf{Fix}: \textbf{Right-Rotate($v$)}.
\end{itemize}

\subsubsection{Case 2: Right-Right (RR)}
\begin{itemize}
    \item \textbf{Condition}: $v$ is heavy on the right ($bal(v) = -2$) AND $x$ is heavy on the right ($bal(x) = -1$ or $0$).
    \item \textbf{Fix}: \textbf{Left-Rotate($v$)}.
\end{itemize}

\subsubsection{Case 3: Left-Right (LR)}
\begin{itemize}
    \item \textbf{Condition}: $v$ is heavy on the left ($bal(v) = +2$) BUT $x$ is heavy on the right ($bal(x) = -1$).
    \item \textbf{Fix}: Double Rotation.
          \begin{enumerate}
              \item \textbf{Left-Rotate($x$)} (Converts structure to LL case).
              \item \textbf{Right-Rotate($v$)}.
          \end{enumerate}
\end{itemize}

\subsubsection{Case 4: Right-Left (RL)}
\begin{itemize}
    \item \textbf{Condition}: $v$ is heavy on the right ($bal(v) = -2$) BUT $x$ is heavy on the left ($bal(x) = +1$).
    \item \textbf{Fix}: Double Rotation.
          \begin{enumerate}
              \item \textbf{Right-Rotate($x$)} (Converts structure to RR case).
              \item \textbf{Left-Rotate($v$)}.
          \end{enumerate}
\end{itemize}

\noindent \textbf{Runtime}: Insertion takes $O(\lg n)$ for the search. Rebalancing requires retracing the path ($O(\lg n)$) and at most \textbf{one} or \textbf{two} rotations (since fixing the lowest imbalance restores the global height). Total: $O(\lg n)$.

\subsection{Deletion}

\textbf{Goal}: Delete node $z$ and restore AVL property.
\textbf{Steps}:
\begin{enumerate}
    \item Perform standard BST deletion.
    \item Retrace the path from the parent of the deleted node (or spliced node) up to the root.
    \item At each node, update height and check balance.
    \item If unbalanced, apply rotations (same 4 cases as insertion).
\end{enumerate}

\subsubsection{Propagation of Imbalance}
Unlike insertion, where one rotation fixes the tree, a rotation during deletion might reduce the height of the subtree, causing the \textit{parent} to become unbalanced.
Thus, rebalancing may propagate all the way up to the root, requiring $O(\lg n)$ rotations.

\noindent \textbf{Runtime}: $O(\lg n)$.

\subsection{Summary}
\begin{table}[h]
    \centering
    \begin{tabular}{|l|c|c|}
        \hline
        \textbf{Operation} & \textbf{BST (Worst Case)} & \textbf{AVL (Worst Case)} \\ \hline
        Search             & $O(n)$                    & $O(\lg n)$                \\ \hline
        Insert             & $O(n)$                    & $O(\lg n)$                \\ \hline
        Delete             & $O(n)$                    & $O(\lg n)$                \\ \hline
    \end{tabular}
    \caption{Comparison of BST and AVL complexities}
\end{table}

\clearpage


\section{Dynamic Programming}

\subsection{The Paradigm Shift}

\subsubsection{Motivation: Divide-and-Conquer vs. Dynamic Programming}
Divide-and-Conquer algorithms (like MergeSort or QuickSort) rely on splitting a problem into \textit{disjoint} subproblems.
\begin{itemize}
    \item \textbf{Disjoint}: The subproblems do not share any common history. Solving the left half of an array tells you nothing about the right half.
    \item \textbf{Efficiency}: Because they are disjoint, we just solve each once and combine them.
\end{itemize}

However, many optimization problems break down into \textbf{Overlapping Subproblems}.
\begin{itemize}
    \item If we use standard recursion, we end up solving the same small subproblems millions of times.
    \item \textbf{Dynamic Programming (DP)} is the technique of "Store, Don't Recompute". We trade space (memory to store results) for time (CPU cycles).
\end{itemize}

\subsection{Motivating Example: Fibonacci Numbers}

The Fibonacci sequence is defined as:
\[ F_0 = 0, F_1 = 1, \quad F_n = F_{n-1} + F_{n-2} \text{ for } n \ge 2 \]

\subsubsection{The Naive Recursive Approach}
\begin{algorithm}
    \caption{Naive-Fib(n)}
    \begin{algorithmic}[1]
        \Procedure{Fib}{$n$}
        \If{$n \le 1$} \Return $n$
        \Else \ \Return \Call{Fib}{$n-1$} + \Call{Fib}{$n-2$}
        \EndIf
        \EndProcedure
    \end{algorithmic}
\end{algorithm}
\textbf{Analysis}:
The recursion tree grows exponentially. To compute $F(5)$, we compute $F(3)$ twice. To compute $F(6)$, we compute $F(4)$ twice and $F(3)$ three times.
\[ T(n) \approx \phi^n \approx 1.618^n \quad (\text{Exponential Explosion}) \]

\subsubsection{The DP Approach (Memoization)}
We use a table to cache results. This is called **Memoization** (Top-Down strategy).
\begin{itemize}
    \item Before computing, check the table.
    \item After computing, save to the table.
\end{itemize}

\begin{algorithm}
    \caption{Memoized-Fib(n)}
    \begin{algorithmic}[1]
        \State Let $memo[0 \dots n]$ be array initialized to $\infty$
        \Procedure{Fib-Memo}{$n$}
        \If{$memo[n] \ne \infty$} \Return $memo[n]$ \Comment{Cache Hit}
        \EndIf
        \If{$n \le 1$} $f = n$
        \Else \ $f = \Call{Fib-Memo}{n-1} + \Call{Fib-Memo}{n-2}$
        \EndIf
        \State $memo[n] = f$ \Comment{Cache Miss: Store result}
        \Return $f$
        \EndProcedure
    \end{algorithmic}
\end{algorithm}
\textbf{Analysis}: Since each $F(i)$ for $0 \le i \le n$ is computed exactly once, the time complexity drops to:
\[ T(n) = \Theta(n) \quad (\text{Linear Time}) \]

\subsection{Case Study: Rod Cutting Problem}

\subsubsection{Problem Definition}
We are given a steel rod of length $n$ and a price table $p_i$ for rod pieces of length $i$. We want to cut the rod into pieces to maximize total revenue $r_n$.

\textbf{Mathematical Formulation}:
We can cut a piece of length $i$ ($1 \le i \le n$) off the left end, and then optimally solve the problem for the remaining length $n-i$.
\[ r_n = \max_{1 \le i \le n} (p_i + r_{n-i}) \]
(Base case: $r_0 = 0$).

\subsubsection{Approach 1: Top-Down with Memoization}
We stick to the recursive logic but add a "memory" (array $r$).
\begin{itemize}
    \item \textbf{Check}: Before computing, is $r[n] \ge 0$? If yes, return it.
    \item \textbf{Save}: After computing max $q$, save $r[n] = q$.
\end{itemize}

\begin{algorithm}
    \caption{Memoized-Cut-Rod(p, n)}\label{alg:memoCutRod}
    \begin{algorithmic}[1]
        \Procedure{Memoized-Cut-Rod}{$p, n$}
        \State Let $r[0 \dots n]$ be new array initialized to $-\infty$
        \State \Return \Call{Memoized-Cut-Rod-Aux}{$p, n, r$}
        \EndProcedure
        \Statex \Comment{Visual separator between procedures}
        \Procedure{Memoized-Cut-Rod-Aux}{$p, n, r$}
        \If{$r[n] \ge 0$}
        \State \Return $r[n]$ \Comment{Return cached value}
        \EndIf
        \If{$n == 0$}
        \State $q = 0$
        \Else
        \State $q = -\infty$
        \For{$i = 1$ \textbf{to} $n$} \Comment{Try every possible first cut}
        \State $q = \max(q, p[i] + \Call{Memoized-Cut-Rod-Aux}{p, n-i, r})$
        \EndFor
        \EndIf
        \State $r[n] = q$ \Comment{Store optimal value for length n}
        \State \Return $q$
        \EndProcedure
    \end{algorithmic}
\end{algorithm}

\subsubsection{Approach 2: Bottom-Up (Tabulation)}
Recursion has overhead (stack depth). We can eliminate recursion by filling the table in order of dependency: from smallest ($0$) to largest ($n$).

\begin{algorithm}
    \caption{Bottom-Up-Cut-Rod(p, n)}\label{alg:bottomUpCutRod}
    \begin{algorithmic}[1]
        \Procedure{Bottom-Up-Cut-Rod}{$p, n$}
        \State Let $r[0 \dots n]$ be new array
        \State $r[0] = 0$
        \For{$j = 1$ \textbf{to} $n$} \Comment{Solve for rod length j}
        \State $q = -\infty$
        \For{$i = 1$ \textbf{to} $j$} \Comment{Try cut length i}
        \State \Comment{Use previously computed solution for remainder j-i}
        \State $q = \max(q, p[i] + r[j-i])$
        \EndFor
        \State $r[j] = q$
        \EndFor
        \Return $r[n]$
        \EndProcedure
    \end{algorithmic}
\end{algorithm}
\textbf{Runtime}: The nested loop structure gives a clear arithmetic sum:
\[ T(n) = \sum_{j=1}^n j = \frac{n(n+1)}{2} = \Theta(n^2) \]

\subsection{Reconstructing the Solution}

Merely knowing the max revenue ($r_n$) is often insufficient; we need the list of cut lengths. To do this, we store the choice that yielded the optimal value.
We extend the algorithm to maintain an array $s[j]$, which stores the \textbf{optimal first cut length} for a rod of length $j$.

\begin{algorithm}
    \caption{Extended-Bottom-Up-Cut-Rod(p, n)}\label{alg:extendedCutRod}
    \begin{algorithmic}[1]
        \Procedure{Extended-Bottom-Up-Cut-Rod}{$p, n$}
        \State Let $r[0 \dots n]$ and $s[0 \dots n]$ be new arrays
        \State $r[0] = 0$
        \For{$j = 1$ \textbf{to} $n$}
        \State $q = -\infty$
        \For{$i = 1$ \textbf{to} $j$}
        \If{$q < p[i] + r[j-i]$}
        \State $q = p[i] + r[j-i]$
        \State $s[j] = i$ \Comment{Record the best cut i}
        \EndIf
        \EndFor
        \State $r[j] = q$
        \EndFor
        \Return $r$ and $s$
        \EndProcedure
    \end{algorithmic}
\end{algorithm}

\subsubsection{Printing the Cuts}
We do not need recursion to print the solution. We simply trace the $s$ array.

\begin{algorithm}
    \caption{Print-Cut-Rod-Solution(p, n)}
    \begin{algorithmic}[1]
        \Procedure{Print-Cut-Rod-Solution}{$p, n$}
        \State $(r, s) = \Call{Extended-Bottom-Up-Cut-Rod}{p, n}$
        \While{$n > 0$}
        \State Print $s[n]$ \Comment{Print first cut}
        \State $n = n - s[n]$ \Comment{Move to remaining length}
        \EndWhile
        \EndProcedure
    \end{algorithmic}
\end{algorithm}

\subsection{Theoretical Foundations}

When should we use DP? Two key properties must hold:

\subsubsection{Optimal Substructure}
A problem exhibits optimal substructure if \textbf{an optimal solution to the problem contains within it optimal solutions to subproblems}.
\begin{itemize}
    \item \textbf{Example}: Shortest Path. If $A \to B \to C$ is the shortest path from $A$ to $C$, then $A \to B$ must be the shortest path from $A$ to $B$.
    \item \textbf{Non-Example}: Longest Simple Path. Longest path $q \to r \to t$ does NOT imply $q \to r$ is the longest simple path.
\end{itemize}
This property justifies the recurrence relation: $r_n = \max(p_i + r_{n-i})$.

\subsubsection{Overlapping Subproblems}
The recursive algorithm visits the \textbf{same} problem instances repeatedly.
\begin{itemize}
    \item This distinguishes DP from Divide-and-Conquer.
    \item \textbf{Merge Sort}: Computes merge on disjoint arrays.
    \item \textbf{Rod Cutting}: To solve for length 4, we need length 3, 2, 1, 0. To solve for length 3, we need 2, 1, 0. Subproblem "length 2" is reused multiple times.
\end{itemize}

\subsection{Summary Comparison}

\begin{table}[h]
    \centering
    \begin{tabular}{|l|l|l|}
        \hline
        \textbf{Feature}     & \textbf{Divide and Conquer}       & \textbf{Dynamic Programming}                   \\ \hline
        \textbf{Subproblems} & Independent (Disjoint)            & Overlapping                                    \\ \hline
        \textbf{Approach}    & Recursive                         & Recursive (Memoized) or Iterative (Tabulation) \\ \hline
        \textbf{Efficiency}  & Splits problem size (e.g., $n/2$) & Reuses solutions (Polynomial time)             \\ \hline
        \textbf{Examples}    & Merge Sort, Quick Sort            & Rod Cutting, Fibonacci, LCS                    \\ \hline
    \end{tabular}
    \caption{Paradigm Comparison}
\end{table}

\clearpage


\section{Greedy Algorithms}

\subsection{The Philosophy of Greed}

\subsubsection{Core Concept}
A \textbf{Greedy Algorithm} builds a solution piece by piece, always choosing the next piece that offers the most immediate benefit.
\begin{itemize}
    \item \textbf{Motto}: "Take what you can get now!"
    \item \textbf{Strategy}: At each step, make a \textbf{locally optimal} choice in the hope that these choices will lead to a \textbf{globally optimal} solution.
    \item \textbf{Benefit}: Usually very efficient (often $O(n \log n)$ or $O(n)$).
    \item \textbf{Risk}: It doesn't always work. It never looks back to correct previous mistakes.
\end{itemize}

\subsubsection{Comparison with Dynamic Programming}
\begin{itemize}
    \item \textbf{Dynamic Programming}: "I need to know the solution to all subproblems before I make a decision." (Bottom-Up or Recursion). It considers \textit{all} possibilities.
    \item \textbf{Greedy}: "I will make a decision right now based on current information, and then I will solve the remaining subproblem." (Top-Down). It commits to a choice \textit{before} solving the subproblem.
\end{itemize}

\subsection{Case Study: Activity Selection Problem}

\subsubsection{The Problem}
We have a set of activities $S = \{a_1, a_2, \dots, a_n\}$. Each activity $a_i$ has a start time $s_i$ and a finish time $f_i$.
\begin{itemize}
    \item \textbf{Constraint}: Two activities are compatible if their time intervals $[s_i, f_i)$ and $[s_j, f_j)$ do not overlap.
    \item \textbf{Goal}: Select the \textbf{maximum number} of mutually compatible activities.
\end{itemize}

\subsubsection{The Greedy Strategy}
How do we choose the "best" activity to pick first?
\begin{enumerate}
    \item \textit{Shortest duration?} No, a short task might happen right in the middle, blocking two other tasks.
    \item \textit{Earliest start time?} No, a task starting at 1 AM might last until midnight, blocking everything else.
    \item \textbf{Earliest Finish Time}: Yes!
\end{enumerate}

\textbf{Intuition}: By picking the activity that finishes \textbf{soonest}, we leave the \textbf{maximum amount of remaining time} for other activities.

\begin{algorithm}
    \caption{Greedy-Activity-Selector(s, f)}\label{alg:activity}
    \begin{algorithmic}[1]
        \Procedure{Greedy-Activity-Selector}{$s, f$}
        \State Sort activities by finish time: $f_1 \le f_2 \le \dots \le f_n$
        \State $A = \{a_1\}$ \Comment{Always select the first activity}
        \State $k = 1$ \Comment{Index of the last selected activity}
        \For{$m = 2$ \textbf{to} $n$}
        \If{$s_m \ge f_k$} \Comment{If activity m starts after k finishes}
        \State $A = A \cup \{a_m\}$
        \State $k = m$
        \EndIf
        \EndFor
        \Return $A$
        \EndProcedure
    \end{algorithmic}
\end{algorithm}

\subsubsection{Runtime Analysis}
\begin{itemize}
    \item Sorting: $O(n \log n)$.
    \item Linear Scan: $\Theta(n)$.
    \item \textbf{Total}: $\Theta(n \log n)$.
\end{itemize}

\subsection{Theoretical Foundations}

To prove a Greedy Algorithm works, the problem must exhibit two key properties:

\subsubsection{Greedy Choice Property}
We can assemble a globally optimal solution by making a locally optimal (greedy) choice.
\begin{itemize}
    \item \textbf{Implication}: We don't need to consider all options (like DP). We just prove that the greedy choice is \textit{always} part of \textit{some} optimal solution.
\end{itemize}

\subsubsection{Optimal Substructure}
An optimal solution to the problem contains within it optimal solutions to subproblems.
\begin{itemize}
    \item If we pick activity $a_1$, the problem reduces to finding optimal activities in the remaining time after $a_1$ finishes.
\end{itemize}

\subsection{The Tale of Two Knapsacks}

This is the classic example showing when Greedy works and when it fails.

\subsubsection{Problem Definitions}
A thief robbing a store finds $n$ items. The $i$-th item has value $v_i$ and weight $w_i$. The knapsack capacity is $W$.
\begin{itemize}
    \item \textbf{0-1 Knapsack}: You must take the item whole or leave it ($x_i \in \{0, 1\}$). (e.g., Gold ingots).
    \item \textbf{Fractional Knapsack}: You can take any fraction of an item ($0 \le x_i \le 1$). (e.g., Gold dust).
\end{itemize}

\subsubsection{The Greedy Strategy: Value Density}
The intuitive greedy choice is to pick items with the highest \textbf{value per unit weight} ($v_i / w_i$).

\subsubsection{Analysis: Why Greedy Fails for 0-1 Knapsack}
Consider Capacity $W=50$.
\begin{itemize}
    \item Item A: Weight 10, Value \$60 (\$6/lb)
    \item Item B: Weight 20, Value \$100 (\$5/lb)
    \item Item C: Weight 30, Value \$120 (\$4/lb)
\end{itemize}
\textbf{Greedy Choice}:
1. Pick Item A (Highest density). Remaining Capacity: 40.
2. Next best is Item B. Remaining Capacity: 20.
3. Item C (Weight 30) doesn't fit.
\textbf{Total Value}: \$160.

\textbf{Optimal Solution}:
Skip Item A. Pick Item B and Item C.
Total Weight: $20+30=50$. \textbf{Total Value}: $\$100 + \$120 = \$220$.

\textbf{Conclusion}: Greedy failed because picking A "polluted" the capacity, preventing us from taking the heavier but valuable combination of B+C. 0-1 Knapsack requires **Dynamic Programming**.



\subsubsection{Analysis: Why Greedy Works for Fractional Knapsack}
Using the same items:
1. Pick Item A (10 lbs). Value \$60. Rem Cap: 40.
2. Pick Item B (20 lbs). Value \$100. Rem Cap: 20.
3. Pick $2/3$ of Item C (20 lbs). Value $(2/3) \times 120 = \$80$.
\textbf{Total Value}: \$240.

\textbf{Conclusion}: Because we can fill the "gaps" with fractions, the "Value Density" strategy guarantees optimality.

\subsection{5. Summary Comparison}

\begin{table}[h]
    \centering
    \begin{tabular}{|l|l|l|}
        \hline
        \textbf{Feature}      & \textbf{Dynamic Programming}                 & \textbf{Greedy Algorithms}               \\ \hline
        \textbf{Choice Logic} & Dependent on subproblems (Look ahead)        & Independent, Local (Look now)            \\ \hline
        \textbf{Structure}    & Bottom-Up or Top-Down                        & Top-Down                                 \\ \hline
        \textbf{Efficiency}   & Slower (often $O(n^2)$ or pseudo-polynomial) & Very Fast (often $O(n \log n)$)          \\ \hline
        \textbf{Correctness}  & Guarantees Optimal                           & Hard to prove (requires Greedy Property) \\ \hline
        \textbf{Example}      & 0-1 Knapsack, LCS, Rod Cutting               & Fractional Knapsack, Activity Selection  \\ \hline
    \end{tabular}
    \caption{DP vs. Greedy}
\end{table}

\clearpage


\section{Elementary Graph Algorithms}

\subsection{Graph Representations}

A graph $G = (V, E)$ can be represented in two standard ways: adjacency lists and adjacency matrices.

\subsubsection{Adjacency Lists}
An array $Adj$ of $|V|$ lists, one for each vertex. The list $Adj[u]$ contains all vertices $v$ such that $(u, v) \in E$.
\begin{itemize}
    \item \textbf{Space Complexity}: $\Theta(V + E)$.
    \item \textbf{Pros}: Compact for sparse graphs.
    \item \textbf{Cons}: Checking if an edge $(u, v)$ exists takes $O(deg(u))$ time.
\end{itemize}

\subsubsection{Adjacency Matrix}
A $|V| \times |V|$ matrix $A$ where $A_{ij} = 1$ if $(i, j) \in E$, and 0 otherwise.
\begin{itemize}
    \item \textbf{Space Complexity}: $\Theta(V^2)$.
    \item \textbf{Pros}: Checking edge existence takes $O(1)$.
    \item \textbf{Cons}: Wasteful for sparse graphs.
\end{itemize}


\subsection{Breadth-First Search (BFS)}

BFS explores a graph by visiting all neighbors of a node before moving to the next level of neighbors. It computes the shortest path distance (in terms of number of edges) from a source vertex $s$ to all reachable vertices.

\subsubsection{Algorithm}
BFS uses a \textbf{queue} to manage the frontier of exploration. Vertices are colored to track their status:
\begin{itemize}
    \item \textbf{White}: Undiscovered.
    \item \textbf{Gray}: Discovered but not fully processed (in the queue).
    \item \textbf{Black}: Fully processed (dequeued).
\end{itemize}

\begin{algorithm}
    \caption{Breadth-First Search}\label{alg:bfs}
    \begin{algorithmic}[1]
        \Procedure{BFS}{$G, s$}
        \For{each vertex $u \in G.V - \{s\}$}
        \State $u.color = \text{WHITE}$
        \State $u.d = \infty$
        \State $u.\pi = \text{NIL}$
        \EndFor
        \State $s.color = \text{GRAY}$
        \State $s.d = 0$
        \State $s.\pi = \text{NIL}$
        \State $Q = \emptyset$
        \State \Call{Enqueue}{$Q, s$}
        \While{$Q \ne \emptyset$}
        \State $u = \Call{Dequeue}{Q}$
        \For{each $v \in G.Adj[u]$}
        \If{$v.color == \text{WHITE}$}
        \State $v.color = \text{GRAY}$
        \State $v.d = u.d + 1$
        \State $v.\pi = u$
        \State \Call{Enqueue}{$Q, v$}
        \EndIf
        \EndFor
        \State $u.color = \text{BLACK}$
        \EndWhile
        \EndProcedure
    \end{algorithmic}
\end{algorithm}

\subsubsection{Analysis}
\begin{itemize}
    \item \textbf{Time Complexity}:
          \begin{itemize}
              \item Each vertex is enqueued and dequeued at most once: $O(V)$.
              \item The adjacency list of each vertex is scanned exactly once: $\sum_{u \in V} |Adj[u]| = \Theta(E)$.
          \end{itemize}
          Total runtime: $O(V + E)$.
    \item \textbf{Shortest Paths}: Upon termination, $v.d = \delta(s, v)$ (shortest path distance) for all reachable $v$.
    \item \textbf{BFS Tree}: The predecessor pointers $\pi$ define a breadth-first tree rooted at $s$.
\end{itemize}

\subsection{Depth-First Search (DFS)}

DFS explores edges out of the most recently discovered vertex $v$ that still has unexplored edges leaving it. Once all edges from $v$ have been explored, the search "backtracks" to explore edges leaving the vertex from which $v$ was discovered.

\subsubsection{Algorithm}
DFS uses recursion (implicit stack). It records two timestamps for each vertex:
\begin{itemize}
    \item $v.d$: Discovery time (vertex becomes Gray).
    \item $v.f$: Finish time (vertex becomes Black).
\end{itemize}

\begin{algorithm}
    \caption{Depth-First Search}\label{alg:dfs}
    \begin{algorithmic}[1]
        \Procedure{DFS}{$G$}
        \For{each vertex $u \in G.V$}
        \State $u.color = \text{WHITE}$
        \State $u.\pi = \text{NIL}$
        \EndFor
        \State $time = 0$
        \For{each vertex $u \in G.V$}
        \If{$u.color == \text{WHITE}$}
        \State \Call{DFS-Visit}{$G, u$}
        \EndIf
        \EndFor
        \EndProcedure
        \Statex
        \Procedure{DFS-Visit}{$G, u$}
        \State $time = time + 1$
        \State $u.d = time$
        \State $u.color = \text{GRAY}$
        \For{each $v \in G.Adj[u]$}
        \If{$v.color == \text{WHITE}$}
        \State $v.\pi = u$
        \State \Call{DFS-Visit}{$G, v$}
        \EndIf
        \EndFor
        \State $u.color = \text{BLACK}$
        \State $time = time + 1$
        \State $u.f = time$
        \EndProcedure
    \end{algorithmic}
\end{algorithm}

\subsubsection{Analysis}
\begin{itemize}
    \item \textbf{Time Complexity}:
          \begin{itemize}
              \item Initialization loops: $\Theta(V)$.
              \item \texttt{DFS-Visit} is called exactly once for each vertex.
              \item The loop over $Adj[u]$ executes $|Adj[u]|$ times. Total cost for edge exploration is $\Theta(E)$.
          \end{itemize}
          Total runtime: $\Theta(V + E)$.
    \item \textbf{Parenthesis Theorem}: For any two vertices $u$ and $v$, exactly one of the following holds:
          \begin{itemize}
              \item Intervals $[u.d, u.f]$ and $[v.d, v.f]$ are disjoint.
              \item Interval $[u.d, u.f]$ is contained within $[v.d, v.f]$ ($v$ is an ancestor of $u$).
              \item Interval $[v.d, v.f]$ is contained within $[u.d, u.f]$ ($u$ is an ancestor of $v$).
          \end{itemize}
\end{itemize}


\subsubsection{Edge Classification}
DFS can classify edges $(u, v)$ based on the color of $v$ when the edge is explored:
\begin{enumerate}
    \item \textbf{Tree Edge}: $v$ is WHITE. (Edge in the DFS forest).
    \item \textbf{Back Edge}: $v$ is GRAY. (Edge to an ancestor). Indicates a cycle.
    \item \textbf{Forward Edge}: $v$ is BLACK and $u.d < v.d$. (Edge to a descendant).
    \item \textbf{Cross Edge}: $v$ is BLACK and $v.d < u.d$. (All other edges).
\end{enumerate}

\subsection{Applications of DFS}

\subsubsection{Topological Sort}
A topological sort of a DAG (Directed Acyclic Graph) is a linear ordering of vertices such that for every edge $(u, v)$, $u$ appears before $v$.
\begin{itemize}
    \item \textbf{Algorithm}: Run DFS. As each vertex is finished (blackened), insert it onto the front of a linked list.
    \item \textbf{Runtime}: $\Theta(V + E)$.
\end{itemize}

\subsubsection{Strongly Connected Components (SCC)}
A strongly connected component of a directed graph is a maximal set of vertices $C \subseteq V$ such that for every pair $u, v \in C$, $u$ is reachable from $v$ and $v$ is reachable from $u$.
\begin{itemize}
    \item \textbf{Algorithm (Kosaraju's)}:
          \begin{enumerate}
              \item Call DFS on $G$ to compute finish times $u.f$.
              \item Compute $G^T$ (transpose of $G$).
              \item Call DFS on $G^T$, but in the main loop consider vertices in order of decreasing $u.f$.
              \item Output the vertices of each tree in the DFS forest of step 3 as a separate SCC.
          \end{enumerate}
    \item \textbf{Runtime}: $\Theta(V + E)$.
\end{itemize}

\clearpage


\section{Depth First Search \& Applications}

\subsection{Review of DFS}

Depth-First Search (DFS) explores edges out of the most recently discovered vertex $v$ that still has unexplored edges leaving it. It uses colors (White, Gray, Black) and timestamps ($v.d$ for discovery, $v.f$ for finish) to track progress.

\begin{algorithm}
    \caption{Depth-First Search}\label{alg:dfsRevisit}
    \begin{algorithmic}[1]
        \Procedure{DFS}{$G$}
        \For{each vertex $u \in G.V$}
        \State $u.color = \text{WHITE}$
        \State $u.\pi = \text{NIL}$
        \EndFor
        \State $time = 0$
        \For{each vertex $u \in G.V$}
        \If{$u.color == \text{WHITE}$}
        \State \Call{DFS-Visit}{$G, u$}
        \EndIf
        \EndFor
        \EndProcedure
    \end{algorithmic}
\end{algorithm}

\begin{algorithm}
    \caption{DFS-Visit}\label{alg:dfsVisit}
    \begin{algorithmic}[1]
        \Procedure{DFS-Visit}{$G, u$}
        \State $time = time + 1$
        \State $u.d = time$
        \State $u.color = \text{GRAY}$
        \For{each $v \in G.Adj[u]$}
        \If{$v.color == \text{WHITE}$}
        \State $v.\pi = u$
        \State \Call{DFS-Visit}{$G, v$}
        \EndIf
        \EndFor
        \State $u.color = \text{BLACK}$
        \State $time = time + 1$
        \State $u.f = time$
        \EndProcedure
    \end{algorithmic}
\end{algorithm}

\textbf{Runtime}: $\Theta(V + E)$.

\subsection{Properties of DFS}

\subsubsection{Parenthesis Structure}
For any two vertices $u$ and $v$, exactly one of the following holds:
\begin{enumerate}
    \item The intervals $[u.d, u.f]$ and $[v.d, v.f]$ are entirely disjoint (neither is a descendant of the other).
    \item The interval $[u.d, u.f]$ is contained entirely within $[v.d, v.f]$ ($u$ is a descendant of $v$).
    \item The interval $[v.d, v.f]$ is contained entirely within $[u.d, u.f]$ ($v$ is a descendant of $u$).
\end{enumerate}

\subsubsection{White-Path Theorem}
\begin{theorem}
    In a depth-first forest of a graph $G=(V,E)$, vertex $v$ is a descendant of vertex $u$ if and only if at the time $u.d$ that the search discovers $u$, there is a path from $u$ to $v$ consisting entirely of white vertices.
\end{theorem}

\subsubsection{Edge Classification}
DFS classifies edges $(u, v)$ based on the color of $v$ when the edge is explored:
\begin{itemize}
    \item \textbf{Tree Edge}: $v$ is WHITE.
    \item \textbf{Back Edge}: $v$ is GRAY. (Implies a cycle).
    \item \textbf{Forward Edge}: $v$ is BLACK and $u.d < v.d$.
    \item \textbf{Cross Edge}: $v$ is BLACK and $u.d > v.d$.
\end{itemize}
\textbf{Note}: In an undirected graph, every edge is either a tree edge or a back edge.

\subsection{Applications}

\subsubsection{Cycle Detection}
\begin{theorem}
    A graph $G$ contains a cycle if and only if a DFS yields at least one back edge.
\end{theorem}
\begin{proof}
    $(\Leftarrow)$ If $(u,v)$ is a back edge, then $v$ is an ancestor of $u$ in the DFS tree. Thus, there is a tree path from $v$ to $u$, and the edge $(u,v)$ completes the cycle.
    $(\Rightarrow)$ If $G$ has a cycle $C$, let $v$ be the first vertex in $C$ discovered by DFS. Let $(u,v)$ be the preceding edge in $C$. At time $v.d$, the vertices of $C$ form a white path from $v$ to $u$. By the White-Path Theorem, $u$ becomes a descendant of $v$. Thus, $(u,v)$ is a back edge.
\end{proof}

\subsubsection{Topological Sort}
A topological sort of a DAG (Directed Acyclic Graph) is a linear ordering of vertices such that for every edge $(u, v)$, $u$ appears before $v$.

\begin{algorithm}
    \caption{Topological Sort}\label{alg:topologicalSort}
    \begin{algorithmic}[1]
        \Procedure{Topological-Sort}{$G$}
        \State Call \Call{DFS}{$G$} to compute finishing times $v.f$ for each vertex $v$
        \State As each vertex is finished, insert it onto the front of a linked list
        \State \Return the linked list of vertices
        \EndProcedure
    \end{algorithmic}
\end{algorithm}

\textbf{Correctness}: If there is an edge $(u,v)$, then $v.f < u.f$.
\begin{itemize}
    \item When $(u,v)$ is explored, $v$ cannot be GRAY (otherwise it's a back edge, implying a cycle).
    \item If $v$ is WHITE, it becomes a descendant of $u$, so $v.f < u.f$.
    \item If $v$ is BLACK, it is already finished, so $v.f < u.d < u.f$.
\end{itemize}
\textbf{Runtime}: $\Theta(V + E)$.

\subsubsection{Strongly Connected Components (SCC)}
An SCC is a maximal set of vertices $C \subseteq V$ such that for every pair $u, v \in C$, $u \leadsto v$ and $v \leadsto u$.

\begin{algorithm}
    \caption{SCC Algorithm (Kosaraju)}\label{alg:scc}
    \begin{algorithmic}[1]
        \Procedure{Strongly-Connected-Components}{$G$}
        \State Call \Call{DFS}{$G$} to compute finishing times $u.f$
        \State Compute $G^T$ (transpose of $G$)
        \State Call \Call{DFS}{$G^T$}, but in the main loop consider vertices in order of decreasing $u.f$ (from step 1)
        \State Output the vertices of each tree in the DFS forest of step 3 as a separate SCC
        \EndProcedure
    \end{algorithmic}
\end{algorithm}

\textbf{Runtime}: $\Theta(V + E)$.
\begin{itemize}
    \item DFS on $G$: $\Theta(V + E)$.
    \item Compute $G^T$: $\Theta(V + E)$.
    \item DFS on $G^T$: $\Theta(V + E)$.
\end{itemize}

\begin{remark}
    This algorithm works because components in the component graph (which is a DAG) are effectively topologically sorted by finish times. Processing nodes in decreasing order of finish times in $G^T$ isolates the SCCs one by one (sink components in the component graph of $G$ become source components in $G^T$).
\end{remark}

\clearpage


\section{Minimum Spanning Trees and Shortest Paths}

\subsection{Minimum Spanning Trees (MST)}

\subsubsection{Problem Definition}

Given a connected, undirected graph $G=(V, E)$ where each edge $(u, v) \in E$ has a weight $w(u, v) \ge 0$. We want to find a subset of edges $T \subseteq E$ that connects all vertices and minimizes the total weight:
\[ w(T) = \sum_{(u, v) \in T} w(u, v) \]
Since $T$ connects all vertices and has minimum weight, it must be acyclic. Thus, $T$ is a tree, called a \textbf{Spanning Tree}.

\subsubsection{Generic Greedy Approach}

We grow a set of edges $A$, maintaining the loop invariant that $A$ is a subset of some minimum spanning tree. At each step, we add a "safe edge" $(u, v)$ to $A$ such that $A \cup \{(u, v)\}$ is still a subset of an MST.

\begin{definition}[Cut]\index{Cut}
    A \textbf{cut} $(S, V-S)$ of an undirected graph $G=(V, E)$ is a partition of $V$.
\end{definition}

\begin{definition}[Cross]\index{Cross}
    An edge $(u, v) \in E$ \textbf{crosses} the cut $(S, V-S)$ if one of its endpoints is in $S$ and the other is in $V-S$.
\end{definition}

\begin{definition}[Light Edge]\index{Light Edge}
    An edge is a \textbf{light edge} crossing a cut if its weight is the minimum of any edge crossing the cut.
\end{definition}

\begin{theorem}[Cut Property / Safe Edge Theorem]\index{Cut Property}
    Let $A$ be a subset of some MST $(S, V-S)$ be any cut that respects $A$ (i.e., no edge in $A$ crosses the cut). Let $(u, v)$ be a light edge crossing $(S, V-S)$. Then, edge $(u, v)$ is safe for $A$.
\end{theorem}

\subsection{Kruskal's Algorithm}

\textbf{Idea}: Kruskal's algorithm builds the MST by finding safe edges from the lightest to the heaviest. It treats the algorithm as managing a forest of trees. Initially, each vertex is its own tree.
\begin{enumerate}
    \item Sort all edges in non-decreasing order of their weights.
    \item Iterate through the sorted edges. For each edge $(u, v)$, if $u$ and $v$ are in different trees (sets), add $(u, v)$ to the forest and combine the two trees.
    \item Repeat until all vertices are in the same tree.
\end{enumerate}

\subsubsection{Disjoint Set Data Structure}
To implement Kruskal's efficiently, we use a \textbf{Disjoint Set Union (DSU)} data structure supporting:
\begin{itemize}
    \item \texttt{Make-Set(v)}: Creates a new set containing only $v$.
    \item \texttt{Find-Set(v)}: Returns a representative element of the set containing $v$.
    \item \texttt{Union(u, v)}: Merges the sets containing $u$ and $v$.
\end{itemize}

\begin{algorithm}
    \caption{Kruskal's Algorithm}\label{alg:kruskal}
    \begin{algorithmic}[1]
        \Procedure{MST-Kruskal}{$G, w$}
        \State $A = \emptyset$
        \For{each vertex $v \in G.V$}
        \State \Call{Make-Set}{$v$}
        \EndFor
        \State Sort the edges of $G.E$ into nondecreasing order by weight $w$
        \For{each edge $(u, v)$ taken from the sorted list}
        \If{\Call{Find-Set}{$u$} $\ne$ \Call{Find-Set}{$v$}}
        \State $A = A \cup \{(u, v)\}$
        \State \Call{Union}{$u, v$}
        \EndIf
        \EndFor
        \State \Return $A$
        \EndProcedure
    \end{algorithmic}
\end{algorithm}

\subsubsection{Analysis}
\begin{itemize}
    \item Initializing sets: $O(V)$.
    \item Sorting edges: $O(E \lg E)$.
    \item Disjoint-set operations: $O(E \alpha(V))$, where $\alpha$ is the inverse Ackermann function (extremely slow-growing, effectively constant).
\end{itemize}
Total runtime: $O(E \lg E)$. Since $E < V^2$, $\lg E = O(\lg V)$, so the time is dominated by sorting: $O(E \lg V)$.

\subsection{Prim's Algorithm}

\textbf{Idea}: Prim's algorithm works like Dijkstra's algorithm. It grows a single tree $T$ from an arbitrary root $r$. At each step, it adds the lightest edge connecting a vertex in $T$ to a vertex outside $T$.

\begin{itemize}
    \item The vertices not yet in the tree reside in a \textbf{Min-Priority Queue} $Q$.
    \item The key of a vertex $v$, $v.key$, is the minimum weight of any edge connecting $v$ to a vertex in the tree.
    \item $v.\pi$ is the parent of $v$ in the tree.
\end{itemize}

\begin{algorithm}
    \caption{Prim's Algorithm}\label{alg:prim}
    \begin{algorithmic}[1]
        \Procedure{MST-Prim}{$G, w, r$}
        \For{each $u \in G.V$}
        \State $u.key = \infty$
        \State $u.\pi = \text{NIL}$
        \EndFor
        \State $r.key = 0$
        \State $Q = G.V$ \Comment{Build Min-Priority Queue}
        \While{$Q \ne \emptyset$}
        \State $u = \Call{Extract-Min}{Q}$
        \For{each $v \in G.Adj[u]$}
        \If{$v \in Q$ \textbf{and} $w(u, v) < v.key$}
        \State $v.\pi = u$
        \State $v.key = w(u, v)$
        \State \Call{Decrease-Key}{$Q, v, w(u, v)$}
        \EndIf
        \EndFor
        \EndWhile
        \EndProcedure
    \end{algorithmic}
\end{algorithm}

\subsubsection{Analysis}
The runtime depends on the implementation of the priority queue:
\begin{itemize}
    \item \textbf{Binary Heap}:
          \begin{itemize}
              \item \texttt{Build-Heap}: $O(V)$.
              \item \texttt{Extract-Min}: $O(\lg V)$, called $V$ times. Total $O(V \lg V)$.
              \item \texttt{Decrease-Key}: $O(\lg V)$, called at most $E$ times. Total $O(E \lg V)$.
          \end{itemize}
          Total: $O(E \lg V)$.
    \item \textbf{Fibonacci Heap}:
          \begin{itemize}
              \item \texttt{Extract-Min}: $O(\lg V)$ amortized.
              \item \texttt{Decrease-Key}: $O(1)$ amortized.
          \end{itemize}
          Total: $O(E + V \lg V)$.
\end{itemize}

\subsection{Single-Source Shortest Paths}

\textbf{Problem}: Given a weighted, directed graph $G=(V, E)$ with weight function $w: E \to \mathbb{R}$, find the shortest path from a source $s$ to all other vertices $v$.
The weight of a path $p = \langle v_0, v_1, \dots, v_k \rangle$ is $w(p) = \sum_{i=1}^k w(v_{i-1}, v_i)$.
The shortest-path weight $\delta(u, v)$ is $\min \{ w(p) : u \leadsto v \}$.

\subsubsection{Relaxation}
The algorithms rely on the technique of \textbf{relaxation}. For each vertex $v$, we maintain an attribute $v.d$, which is an upper bound on the weight of a shortest path from source $s$ to $v$.

\begin{algorithm}
    \caption{Initialization and Relaxation}\label{alg:relax}
    \begin{algorithmic}[1]
        \Procedure{Initialize-Single-Source}{$G, s$}
        \For{each vertex $v \in G.V$}
        \State $v.d = \infty$
        \State $v.\pi = \text{NIL}$
        \EndFor
        \State $s.d = 0$
        \EndProcedure
        \Statex
        \Procedure{Relax}{$u, v, w$}
        \If{$v.d > u.d + w(u, v)$}
        \State $v.d = u.d + w(u, v)$
        \State $v.\pi = u$
        \EndIf
        \EndProcedure
    \end{algorithmic}
\end{algorithm}

\subsection{Dijkstra's Algorithm}

\textbf{Assumption}: Edge weights are non-negative ($w(u, v) \ge 0$).

\textbf{Idea}: Dijkstra's algorithm maintains a set $S$ of vertices whose final shortest-path weights have been determined. It repeatedly selects the vertex $u \in V - S$ with the minimum shortest-path estimate, adds $u$ to $S$, and relaxes all edges leaving $u$.

\begin{algorithm}
    \caption{Dijkstra's Algorithm}\label{alg:dijkstra}
    \begin{algorithmic}[1]
        \Procedure{Dijkstra}{$G, w, s$}
        \State \Call{Initialize-Single-Source}{$G, s$}
        \State $S = \emptyset$
        \State $Q = G.V$ \Comment{Min-priority queue keyed by $d$ values}
        \While{$Q \ne \emptyset$}
        \State $u = \Call{Extract-Min}{Q}$
        \State $S = S \cup \{u\}$
        \For{each vertex $v \in G.Adj[u]$}
        \State \Call{Relax}{$u, v, w$} \Comment{Implicitly performs Decrease-Key if $v \in Q$}
        \EndFor
        \EndWhile
        \EndProcedure
    \end{algorithmic}
\end{algorithm}

\subsubsection{Correctness}
\begin{theorem}
    Dijkstra's algorithm, run on a weighted, directed graph $G$ with non-negative weight function $w$ and source $s$, terminates with $u.d = \delta(s, u)$ for all vertices $u \in V$.
\end{theorem}
\begin{proof}
    By invariant: At the start of each iteration of the \textbf{while} loop, $v.d = \delta(s, v)$ for each vertex $v \in S$.
    When a vertex $u$ is added to $S$, its shortest path is already found. This relies on the fact that since weights are non-negative, taking a path through a vertex with a larger $d$ value cannot result in a shorter path to $u$.
\end{proof}

\subsubsection{Analysis}
\begin{itemize}
    \item \textbf{Using Binary Heap}: Each vertex extracted once: $O(V \lg V)$. Each edge relaxed once: $O(E \lg V)$. Total: $O(E \lg V)$.
    \item \textbf{Using Fibonacci Heap}: $O(V \lg V + E)$.
\end{itemize}

\subsection{Summary of Algorithms}

\begin{table}[h]
    \centering
    \begin{tabular}{|l|l|l|l|}
        \hline
        \textbf{Algorithm} & \textbf{Problem} & \textbf{Constraint}  & \textbf{Runtime (Binary Heap)} \\ \hline
        Kruskal            & MST              & Undirected           & $O(E \lg V)$                   \\ \hline
        Prim               & MST              & Undirected           & $O(E \lg V)$                   \\ \hline
        Dijkstra           & Shortest Path    & Non-negative weights & $O(E \lg V)$                   \\ \hline
        Bellman-Ford       & Shortest Path    & No negative cycles   & $O(VE)$                        \\ \hline
    \end{tabular}
    \caption{Comparison of Graph Algorithms}
\end{table}

\clearpage

\printindex

\end{document}